% संपूर्ण मराठी ग्रंथातील प्रकरण 4: संचांचे आकारमान
% हा संपादनयोग्य स्रोतखंड आहे. संकलनासाठी संपूर्ण स्रोतसंग्रहातील
% अक्षररूपे, पूर्वभाग, चिन्हव्याख्या आणि आकृती-स्रोत आवश्यक आहेत.
% मूळ: Open Logic Project; मजकूर आणि रूपांतर CC BY 4.0.
% यंत्रानुवाद, दुरुस्ती आणि तपासणी: OpenAI Codex —
% GPT-5.6 Sol आणि GPT-6 Sol, दोन्ही Ultra effort.
% दुरुस्ती-पुनरावलोकन: GPT-6.1 Sol, Ultra effort.
\chapter{संचांचे आकारमान}\label{sfr:siz::chap}






\section{परिचय}\label{sfr:siz:int:sec}

1870 च्या दशकात गेओर्क कँटर यांनी संच सिद्धांत
विकसित केला, तेव्हा अनंत समूहाची कल्पना---मध्ययुगीन विचारवंतांनी
म्हटल्याप्रमाणे प्रत्यक्ष अस्तित्वातील अनंतता---स्वीकारार्ह करणे हे
त्यांच्या उद्दिष्टांपैकी एक होते. वेगवेगळ्या संचांच्या \emph{आकारमानाचा}
त्यांनी केलेला विचार हा त्यातील महत्त्वाचा भाग होता. $a$, $b$ आणि $c$
हे सर्व भिन्न असतील, तर $\{a, b, c\}$ हा संच $\{a, b\}$ पेक्षा
\emph{मोठा} आहे असे सहज वाटते. पण अनंत संचांचे काय? ते सर्व
एकमेकांएवढेच मोठे असतात का? तसे नसते, हे दिसून येते.

येथील पहिली महत्त्वाची कल्पना म्हणजे प्रगणन. कोणत्याही सांत संचातील
सर्व घटक लिहून त्याची यादी देता येते. काही अनंत संचांतील
सर्व घटक चीही यादी देता येते, जर यादी स्वतः अनंत असण्याची
मुभा दिली, तर. अशा संचांना गणनीय म्हणतात. काही अनंत संच
गणनीय नसतात, हा कँटर यांचा आश्चर्यकारक निष्कर्ष होता.
या प्रकरणाच्या अखेरीस तो आपल्याला पूर्णपणे समजेल.









\section{प्रगणने आणि गणनीय संच}\label{sfr:siz:enm:sec}

\begin{editorial}
  या विभागात संचांच्या प्रगणनांची चर्चा आहे. त्यांची व्याख्या
  $\PosInt$ पासूनच्या आच्छादनांनी केली आहे. गणिताची फारशी पार्श्वभूमी
  नसलेल्या वाचकांसाठी विषय हळूहळू मांडला आहे. पर्यायी, अधिक संक्षिप्त
  आवृत्ती \ref{sfr:siz:enm-alt:sec} मध्ये दिली आहे. तेथे प्रगणनांची वेगळी
  व्याख्या केली आहे: $\Nat$ किंवा त्याच्या आरंभीच्या खंडाशी
  एकास-एक आच्छादन असणे.
\end{editorial}

\begin{explain}
संचांचे घटक यादीत देऊन आपण आधीच संचांची उदाहरणे दिली आहेत.
संच अनंत असला, तरी त्याच्या घटक ची यादी कशी आणि कधी
देता येते, याची आता अधिक व्यापकपणे चर्चा करू.
\end{explain}

\begin{defn}[प्रगणनाची अनौपचारिक व्याख्या]
अनौपचारिकपणे, संच~$A$ चे \emph{प्रगणन} म्हणजे~$A$ मधील
घटक ची अशी यादी (कदाचित अनंत) की $A$ मधील प्रत्येक
घटक त्या यादीत कोणत्या तरी सांत क्रमांकाच्या स्थानी येतो.
$A$ चे प्रगणन असेल, तर $A$ हा \emph{गणनीय} आहे असे म्हणतात.
\end{defn}

\begin{explain}
प्रगणनांविषयी काही मुद्दे:
\begin{enumerate}
\item ज्या यादीला सुरुवात असते आणि जिच्यात पहिल्या घटकाव्यतिरिक्त
  प्रत्येक घटक च्या लगेच आधी एकच घटक असतो, अशीच यादी
  आपण प्रगणन मानतो. दुसऱ्या शब्दांत, यादीच्या पहिल्या घटक
  आणि इतर कोणत्याही घटक यांच्यामध्ये केवळ सांत संख्येने घटक
  असतात. विशेषतः, प्रगणनातील प्रत्येक घटक चे स्थान सांत
  क्रमांकाचे असते: पहिल्या घटक चे स्थान~$1$, दुसऱ्याचे~$2$, इत्यादी.
\item एकाच संच~$A$ ची वेगवेगळी प्रगणने असू शकतात, ज्यांत
  घटक येण्याचा क्रम वेगळा असतो: $4$, $1$, $25$, $16$,~$9$
  हे पहिल्या पाच वर्गसंख्यांच्या संचाचे प्रगणन आहे, तसेच
  $1$, $4$, $9$, $16$,~$25$ हेही आहे.
\item पुनरुक्ती असलेली प्रगणनेही प्रगणनेच असतात: $1$, $1$, $2$,
  $2$, $3$, $3$,~\dots{} हे त्याच संचाचे प्रगणन आहे, ज्याचे
  प्रगणन $1$, $2$, $3$,~\dots{} हे आहे.
\item विशिष्ट प्रगणन सांगताना क्रम आणि पुनरुक्ती यांना महत्त्व
  \emph{असते}: धन पूर्णांकांचे प्रगणन $1$, $2$, $3$, $1$, \dots{}
  अशा सुरुवातीने करता येते; पण नेहमीच्या $1$, $2$, $3$, $4$,~\dots
  या पद्धतीने प्रगणन केले, तर त्यातील नमुना अधिक सहज दिसतो.
\item प्रगणनाला सुरुवात असली पाहिजे: \dots, $3$, $2$, $1$ हे
  धन पूर्णांकांचे प्रगणन नाही, कारण त्याला पहिला घटक नाही.
  अनौपचारिक व्याख्येवरून हे कसे येते ते पाहण्यासाठी स्वतःला विचारा:
  ``या यादीत 76 ही संख्या कोणत्या स्थानी येते?''
\item पुढील यादी धन पूर्णांकांचे प्रगणन नाही:
  $1$, $3$, $5$, \dots, $2$, $4$, $6$, \dots\@ अडचण अशी की
  सम संख्या सांत क्रमांकांच्या स्थानांवर येण्याऐवजी
  $\infty + 1$, $\infty + 2$, $\infty + 3$ या स्थानांवर येतात.
\item रिक्त संच गणनीय आहे: रिक्त यादी हे त्याचे प्रगणन आहे!{}
\end{enumerate}
\end{explain}

\begin{prop}
  $A$ चे प्रगणन असेल, तर त्याचे पुनरुक्ती नसलेले प्रगणनही असते.
\end{prop}

\begin{proof}
  $A$ चे $x_1$, $x_2$, \dots{} असे प्रगणन आहे, आणि प्रत्येक
  $x_i$ हा~$A$ चा घटक आहे असे समजा. प्रगणनातून पुन्हा
  आलेले घटक काढून टाकून त्यातील पुनरुक्ती काढता येते.
  उदाहरणार्थ, पुढीलप्रमाणे नवे प्रगणन मिळवता येते: $x_i$ हा~$A$
  मधील घटक असून $x_1$, \dots, $x_{i-1}$ यांमध्ये नसेल,
  तर त्याला यादीत लिहायचे; आणि $x_i$ हा आधीच $x_1$, \dots,~$x_{i-1}$
  यांमध्ये असेल, तर त्याला यादीतून काढून टाकायचे.
\end{proof}

मागील युक्तिवाद दाखवतो की प्रगणने आणि गणनीय संच नीट
समजून घेण्यासाठी, तसेच त्यांविषयी गोष्टी सिद्ध करण्यासाठी,
अधिक नेमकी व्याख्या हवी. ती पुढे दिली आहे.

\begin{defn}[प्रगणनाची औपचारिक व्याख्या]
$A \neq \emptyset$ या संचाचे \emph{प्रगणन} म्हणजे कोणतेही
आच्छादक फलन $f \colon \PosInt \to A$.
\end{defn}

\begin{explain}
औपचारिक व्याख्या आणि कदाचित अनंत असणाऱ्या यादीचा वापर करणारी
अनौपचारिक व्याख्या समतुल्य आहेत याची खात्री करू. प्रथम, $\PosInt$
पासून संच~$A$ पर्यंतचे कोणतेही आच्छादक फलन~$A$ चे प्रगणन करते.
अशा फलनाने वरील अनौपचारिक अर्थाचे प्रगणन ठरते: $f(1)$, $f(2)$,
$f(3)$, \dots{} ही यादी. $f$ हे आच्छादक असल्यामुळे~$A$
मधील प्रत्येक घटक हा कोणत्या तरी~$n \in \PosInt$ साठीचे
~$f(n)$ चे मूल्य असतो. म्हणून $A$ मधील प्रत्येक घटक यादीत
सांत क्रमांकाच्या स्थानी येतो. फलन एकास-एक असेलच असे नसल्याने
यादीत पुनरुक्ती असू शकते; पण वर म्हटल्याप्रमाणे ती चालते.

उलट,~$A$ मधील सर्व घटक चे प्रगणन करणारी यादी दिली असल्यास
आच्छादक फलन $f\colon \PosInt \to A$ असे परिभाषित करता येते:
$f(n)$ म्हणजे यादीतील $n$ व्या स्थानी असलेला घटक; आणि
$n$ व्या स्थानी घटक नसेल, तर यादीचा शेवटचा घटक.
$A$ रिक्त असतो आणि म्हणून यादीही रिक्त असते, फक्त या एकाच बाबतीत
या पद्धतीने आच्छादक फलन मिळत नाही. म्हणून प्रत्येक अरिक्त
यादी एक आच्छादक फलन $f\colon \PosInt \to A$ ठरवते.
\end{explain}

\begin{defn}
  \label{sfr:siz:enm:defn:enumerable}
  संच~$A$ हा गणनीय असतो तेव्हा आणि केवळ तेव्हाच,
  जेव्हा तो रिक्त असतो किंवा त्याचे प्रगणन असते.
\end{defn}

\begin{ex}
धन पूर्णांकांचे ($\PosInt$) प्रगणन करणारे फलन म्हणजे फक्त
$f(n) = n$ ने दिलेले एकरूपता फलन. नैसर्गिक संख्या $\Nat$ यांचे
प्रगणन करणारे फलन म्हणजे $g(n) = n - 1$.
\end{ex}

\begin{ex}
$f\colon \PosInt \to \PosInt$ आणि $g \colon \PosInt \to \PosInt$
ही फलने अशी दिली आहेत:
\begin{align*}
f(n) & = 2n \text{ आणि}\\
g(n) & = 2n - 1
\end{align*}
ती अनुक्रमे सम धन पूर्णांकांचे आणि विषम धन पूर्णांकांचे प्रगणन करतात.
मात्र यांपैकी कोणतेही फलन $\PosInt$ चे प्रगणन नाही, कारण कोणतेही
आच्छादक नाही.
\end{ex}

\begin{prob}
धन वर्गसंख्या $1$, $4$, $9$, $16$, \dots{} यांचे प्रगणन परिभाषित करा.
\end{prob}

\begin{ex}
$f(n) = (-1)^{n} \lceil \frac{(n-1)}{2}\rceil$ हे फलन पूर्णांकांचा
संच~$\Int$ प्रगणित करते. येथे $\lceil x \rceil$ हे
\emph{ऊर्ध्व पूर्णांक फलन} दर्शवते; ते $x$ ला त्याच्यापेक्षा कमी
नसलेल्या सर्वांत जवळच्या पूर्णांकाकडे नेते. धन आणि ऋण पूर्णांकांमध्ये
आलटूनपालटून ``उड्या मारून'' $f$ हे $\Int$ ची मूल्ये कशी निर्माण
करते ते पाहा:
\[\begin{array}{c c c c c c c c}
f(1) & f(2) & f(3) & f(4) & f(5) & f(6) & f(7) & \dots \\ \\
- \lceil \tfrac{0}{2} \rceil & \lceil \tfrac{1}{2}\rceil & - \lceil \tfrac{2}{2} \rceil & \lceil \tfrac{3}{2} \rceil & - \lceil \tfrac{4}{2} \rceil  & \lceil \tfrac{5}{2}
\rceil & - \lceil \tfrac{6}{2} \rceil & \dots \\ \\
0 & 1 & -1 & 2 & -2 & 3 & -3 & \dots
\end{array}
\]
\begin{quote}\small\textbf{स्रोतदुरुस्ती.} OLSIZ-001: गोठवलेल्या इंग्रजी तक्त्याच्या शेवटच्या ओळीत f(7) साठी अपेक्षित \ensuremath{-}3 हे मूल्य सुटले आहे; सूत्र आणि वरील दोन ओळी ते मूल्य निश्चित करतात. मराठी तक्त्यात \ensuremath{-}3 समाविष्ट करून ही उणीव दुरुस्त केली आहे.\end{quote}
$f$ ची व्याख्या पुढीलप्रमाणे प्रकरणे पाडून केली आहे असेही मानता येते:
\[f(n) = \begin{cases}
  0 & \text{जर $n = 1$ असेल}\\
  n/2 & \text{जर $n$ सम असेल}\\
  -(n-1)/2 & \text{जर $n$ विषम आणि $>1$ असेल}
  \end{cases}
\]
\end{ex}

\begin{prob}
  $A$ आणि $B$ हे गणनीय असतील, तर $A \cup B$ हाही
  गणनीय असतो, हे दाखवा. यासाठी आच्छादक फलने
  $f\colon \PosInt \to A$ आणि $g\colon \PosInt \to B$ आहेत असे
  समजा. एक आच्छादक फलन~$h\colon \PosInt \to A \cup B$
  परिभाषित करा आणि ते आच्छादक आहे हे सिद्ध करा. $A$ किंवा
  ~$B = \emptyset$ असण्याच्याही बाबतींचा विचार करा.
\end{prob}

\begin{prob}
  $B \subseteq A$ आणि $A$ हा गणनीय असेल, तर~$B$ हाही
  गणनीय असतो, हे दाखवा. यासाठी एक आच्छादक फलन
  $f\colon \PosInt \to A$ आहे असे समजा. एक आच्छादक
  फलन~$g\colon \PosInt \to B$ परिभाषित करा आणि ते
  आच्छादक आहे हे सिद्ध करा. $B = \emptyset$ असल्यास काय होते?
\end{prob}

\begin{prob}
  $A_1$, $A_2$, \dots, $A_n$ हे सर्व गणनीय असतील, तर
  $A_1 \cup \dots \cup A_n$ हाही गणनीय असतो, हे
  $n$ वर विगमन करून दाखवा. $A$ आणि~$B$ हे दोन संच
  गणनीय असल्यास~$A \cup B$ हाही गणनीय असतो
  हे तथ्य तुम्ही गृहीत धरू शकता.
\end{prob}

संचातील घटक ची यादी करताना $1$ व्या घटक पासून
मोजायला सुरुवात करणे कदाचित अधिक स्वाभाविक असले, तरी गणितज्ञांना
मोजणीसाठी नैसर्गिक संख्या~$\Nat$ वापरायला आवडतात. ते यादीतील
$0$ व्या, $1$ व्या, $2$ व्या, अशा घटक विषयी बोलतात.
त्याला अनुसरून प्रगणनाची व्याख्या $\Nat$ पासून~$A$ पर्यंतचे
आच्छादक फलन अशी करता येते. अर्थात, या दोन व्याख्या समतुल्य आहेत.

\begin{prop}\label{sfr:siz:enm:prop:enum-shift}
  आच्छादन $f\colon \PosInt \to A$ असते तेव्हा आणि
  केवळ तेव्हाच, जेव्हा आच्छादन $g\colon \Nat \to A$ असते.
\end{prop}

\begin{proof}
  आच्छादन $f\colon \PosInt \to A$ दिले असता, सर्व
  $n \in \Nat$ साठी $g(n) = f(n+1)$ अशी व्याख्या करता येते.
  $g\colon \Nat \to A$ हे आच्छादक आहे हे सहज दिसते.
  उलट, आच्छादन $g\colon \Nat \to A$ दिले असता,
  $f(n) = g(n-1)$ अशी व्याख्या करा.
\end{proof}

यावरून पुढील निष्कर्ष मिळतो:

\begin{cor}\label{sfr:siz:enm:cor:enum-nat}
संच $A$ हा गणनीय असतो तेव्हा आणि केवळ तेव्हाच, जेव्हा
तो रिक्त असतो किंवा आच्छादक फलन $f\colon \Nat \to A$ असते.
\end{cor}

संच~$A$ मधील घटक च्या यादीतून पुनरुक्ती नसलेली यादी
मिळवता येते याची वर चर्चा केली. प्रगणनांबाबतही हे खरे आहे,
पण ते काटेकोरपणे मांडणे आणि सिद्ध करणे थोडे कठीण आहे. कोणतेही
फलन $f\colon \PosInt \to A$ हे सर्व $n \in \PosInt$ साठी
परिभाषित असले पाहिजे. ~$A$ मध्ये फक्त सांत संख्येने घटक
असतील, तर~$\PosInt$ मधील अनंत संख्येने असलेल्या घटक
वर परिभाषित असणारे आणि~$A$ मधील सर्व घटक मूल्ये म्हणून
घेणारे, पण एकच मूल्य दोनदा कधीच न घेणारे फलन असू शकत नाही.
अशा बाबतीत, म्हणजे पुनरुक्ती नसलेली यादी सांत असते तेव्हा,~$f$
साठी वेगळा प्रांत निवडला पाहिजे; त्यात केवळ सांत संख्येने
घटक असले पाहिजेत. पुनरुक्ती नसणे म्हणजे $f$ हे
एकास-एक असले पाहिजे. ते आच्छादक ही असल्यामुळे
आपण कोणता तरी सांत संच $\{1, \dots, n\}$ किंवा $\PosInt$
आणि~$A$ यांच्यामधील एकास-एक आच्छादन शोधत आहोत.

\begin{prop}\label{sfr:siz:enm:prop:enum-bij}
$f\colon \PosInt \to A$ हे आच्छादक असेल (म्हणजे~$A$
चे प्रगणन असेल), तर एकास-एक आच्छादन $g\colon Z \to A$ असते,
जेथे $Z$ हा~$\PosInt$ किंवा कोणत्या तरी~$n \in \PosInt$
साठी $\{1, \dots, n\}$ असतो.
\end{prop}

\begin{proof}
  फलन $g$ ची व्याख्या आपण पुनरावर्ती रीतीने करतो: $g(1) = f(1)$
  असे ठेवा. $g(i)$ आधीच परिभाषित झाले असेल, तर $f(1)$, $f(2)$,
  \dots{} यांपैकी $g(1)$, \dots, $g(i)$ यांत आधीच न आलेले
  पहिले मूल्य, असे मूल्य असेल तर, $g(i+1)$ म्हणून घ्या.
  $A$ मध्ये नेमके $n$ घटक असतील, तर $g(1)$, \dots, $g(n)$
  ही सर्व मूल्ये परिभाषित असतात, आणि आपल्याला
  $g\colon \{1, \dots, n\} \to A$ असे फलन मिळते.
  $A$ मध्ये अनंत संख्येने घटक असतील, तर कोणत्याही $i$
  साठी $f(1)$, $f(2)$, \dots{} या प्रगणनात~$A$ मधील असा
  घटक असलाच पाहिजे, जो $g(1)$, \dots, $g(i)$ यांत
  आधीच आलेला नाही. अशा बाबतीत आपण $g\colon \PosInt \to A$
  हे फलन परिभाषित केले आहे.

  फलन $g$ हे आच्छादक आहे, कारण~$A$ मधील कोणताही घटक
  $f(1)$, $f(2)$, \dots{} यांमध्ये येतो ($f$ हे आच्छादक
  असल्यामुळे), आणि म्हणून कधी ना कधी तो कोणत्या तरी~$i$
  साठीचे~$g(i)$ चे मूल्य होईल. ते एकास-एक ही आहे:
  $g(j) = g(i)$ असणारे $j < i$ असते, तर $g(i)$ हे मूल्य
  आधीच $g(1)$, \dots, $g(i-1)$ यांमध्ये आलेले असते;
  पण हे~$g$ च्या आपल्या व्याख्येविरुद्ध आहे.
\end{proof}

\begin{cor}\label{sfr:siz:enm:cor:enum-nat-bij}
संच $A$ हा गणनीय असतो तेव्हा आणि केवळ तेव्हाच, जेव्हा
तो रिक्त असतो किंवा एकास-एक आच्छादन $f\colon N \to A$ असते,
जेथे $N = \Nat$ किंवा कोणत्या तरी $n \in \Nat$ साठी
$N = \{0, \dots, n\}$ असते.
\end{cor}

\begin{proof}
$A$ हा गणनीय असतो तेव्हा आणि केवळ तेव्हाच, जेव्हा $A$
रिक्त असतो किंवा आच्छादक $f\colon \PosInt \to A$ असते.
\ref{sfr:siz:enm:prop:enum-bij} नुसार दुसरी अट तेव्हा आणि केवळ तेव्हाच
खरी असते, जेव्हा एकास-एक व आच्छादक फलन~$f\colon Z \to A$ असते,
जेथे $Z = \PosInt$ किंवा कोणत्या तरी $n \in \PosInt$ साठी
$Z = \{1, \dots, n\}$ असते. \ref{sfr:siz:enm:prop:enum-shift} च्या
सिद्धतेतील युक्तिवादानुसार, हे तेव्हा आणि केवळ तेव्हाच शक्य होते,
जेव्हा एकास-एक आच्छादन $g\colon N \to A$ असते, जेथे
$N = \Nat$ किंवा $N = \{0, \dots, n-1\}$ असते.
\begin{quote}\small\textbf{स्रोतदुरुस्ती.} MRSIZ-001: विधानातील सांत प्रांत 0 ते n आणि सिद्धतेतील 0 ते n\ensuremath{-}1 हे वेगवेगळ्या चलनामांनी त्याच सांत आरंभीच्या खंडांच्या कुटुंबाचे वर्णन करतात; येथे स्रोताचे दोन्ही निर्देशांक जसेच्या तसे ठेवले आहेत.\end{quote}
\end{proof}

\begin{prob}
  \ref{sfr:siz:enm:defn:enumerable} नुसार संच $A$ हा गणनीय
  असतो तेव्हा आणि केवळ तेव्हाच, जेव्हा $A = \emptyset$ किंवा
  आच्छादक $f\colon \PosInt \to A$ असते.
  ``गणनीय संच'' याची नेमकी व्याख्या पुढीलप्रमाणेही करता येते:
  संच गणनीय असतो तेव्हा आणि केवळ तेव्हाच, जेव्हा एकास-एक
  फलन $g\colon A \to \PosInt$ असते. या व्याख्या समतुल्य आहेत
  हे दाखवा. म्हणजे, एकास-एक फलन $g\colon A \to \PosInt$
  असते तेव्हा आणि केवळ तेव्हाच, जेव्हा $A = \emptyset$ किंवा
  आच्छादक $f\colon \PosInt \to A$ असते, हे दाखवा.
\end{prob}









\section{कँटर यांची नागमोडी पद्धत}\label{sfr:siz:zigzag:sec}

\begin{explain}
काही ``सोप्या'' प्रगणनांचा आपण आधीच विचार केला आहे. आता जरा कठीण
उदाहरण पाहू. नैसर्गिक संख्यांच्या जोड्यांचा संच विचारात घ्या
; त्याची व्याख्या आपण
\ref{sfr:set:pai:sec} मध्ये पुढीलप्रमाणे केली:
\[\Nat \times \Nat = \Setabs{\tuple{n,m}}{n,m \in \Nat}
\]
या क्रमित जोड्या आपण पुढीलप्रमाणे एका \emph{सरणीत} मांडू शकतो:
\[\begin{array}{ c | c | c | c | c | c}
& \mathbf 0 & \mathbf 1 & \mathbf 2 & \mathbf 3 & \dots \\
\hline
\mathbf 0 & \tuple{0,0} & \tuple{0,1} & \tuple{0,2} & \tuple{0,3} & \dots \\
\hline
\mathbf 1 & \tuple{1,0} & \tuple{1,1} & \tuple{1,2} & \tuple{1,3} & \dots \\
\hline
\mathbf 2 & \tuple{2,0} & \tuple{2,1} & \tuple{2,2} & \tuple{2,3} & \dots \\
\hline
\mathbf 3 & \tuple{3,0} & \tuple{3,1} & \tuple{3,2} & \tuple{3,3} & \dots \\
\hline
\vdots & \vdots & \vdots & \vdots & \vdots & \ddots\\
\end{array}
\]
$\Nat \times \Nat$ मधील प्रत्येक क्रमित जोडी या सरणीत नेमकी एकदाच
येईल, हे स्पष्ट आहे. विशेषतः, $\tuple{n,m}$ ही जोडी $n$ व्या ओळीत
आणि $m$ व्या स्तंभात येईल. पण अशा सरणीतील घटकांची ``एकमितीय'' यादी
कशी करायची? पुढील सरणीतील नमुना ती करण्याचा एक मार्ग दाखवतो
(अर्थातच इतर अनेक मार्गही आहेत):
\[\begin{array}{ c | c | c | c | c | c | c}
& \mathbf 0 & \mathbf 1 & \mathbf 2 & \mathbf 3 & \mathbf 4 &\dots \\
\hline
\mathbf 0 & 0  & 1& 3 & 6& 10 &\ldots \\
\hline
\mathbf 1 &2 & 4& 7 & 11 & \dots &\ldots \\
\hline
\mathbf 2 & 5 & 8 & 12 & \ldots & \dots&\ldots \\
\hline
\mathbf 3 & 9 & 13 & \ldots & \ldots & \dots & \ldots \\
\hline
\mathbf 4 & 14 & \ldots & \ldots & \ldots & \dots & \ldots \\
\hline
\vdots & \vdots & \vdots & \vdots & \vdots&\ldots & \ddots\\
\end{array}
\]\noindent
या नमुन्याला \emph{कँटर यांची नागमोडी पद्धत} म्हणतात. तो
$\Nat \times \Nat$ चे पुढीलप्रमाणे प्रगणन करतो:
\[\tuple{0,0}, \tuple{0,1}, \tuple{1,0}, \tuple{0,2}, \tuple{1,1},
\tuple{2,0}, \tuple{0,3}, \tuple{1,2}, \tuple{2,1}, \tuple{3,0}, \dots
\]
यावरून पुढील विधान सिद्ध होते:
\end{explain}

\begin{prop}\label{sfr:siz:zigzag:natsquaredenumerable}
$\Nat \times \Nat$ हा गणनीय आहे.
\end{prop}

\begin{proof}
प्रत्येक $k \in \Nat$ ला कँटर यांच्या नागमोडी सरणीतील ज्या $n$ व्या
ओळीच्या आणि $m$ व्या स्तंभाच्या जागेचे मूल्य $k$ आहे, त्या
$\tuple{n,m} \in \Nat \times \Nat$ या क्रमित घटकाशी जोडणारे
$f \colon \Nat \to \Nat\times\Nat$ असे फलन घ्या.
\end{proof}

\begin{explain}
ही पद्धत अधिक व्यापक बाबतीतही सहज लागू होते. उदाहरणार्थ, नैसर्गिक
संख्यांच्या क्रमित तीन-घटकसमूहांचा पुढील संच प्रगणित करण्यासाठी ती
वापरता येते:
\[\Nat \times \Nat \times \Nat = \Setabs{\tuple{n,m,k}}{n,m,k \in \Nat}
\]
$\Nat \times \Nat \times \Nat$ हा $\Nat \times \Nat$ आणि $\Nat$
यांचा कार्तीय गुणाकार आहे असे आपण मानतो; म्हणजे,
\[\Nat^3 = (\Nat \times \Nat) \times \Nat =
\Setabs{\tuple{\tuple{n,m},k}}{n, m, k
  \in \Nat }
\]
म्हणून एका अक्षाला $\Nat$ चे प्रगणन आणि दुसऱ्या अक्षाला $\Nat^2$
चे प्रगणन अशी नावे देऊन, एका सरणीच्या साहाय्याने $\Nat^3$ चे
प्रगणन करता येते:
\[\begin{array}{ c | c | c | c | c | c}
& \mathbf 0 & \mathbf 1 & \mathbf 2 & \mathbf 3 & \dots \\
\hline
\mathbf{\tuple{0,0}} & \tuple{0,0,0} & \tuple{0,0,1} & \tuple{0,0,2} & \tuple{0,0,3} & \dots \\
\hline
\mathbf{\tuple{0,1}} & \tuple{0,1,0} & \tuple{0,1,1} & \tuple{0,1,2} & \tuple{0,1,3} & \dots \\
\hline
\mathbf{\tuple{1,0}} & \tuple{1,0,0} & \tuple{1,0,1} & \tuple{1,0,2} & \tuple{1,0,3} & \dots \\
\hline
\mathbf{\tuple{0,2}} & \tuple{0,2,0} & \tuple{0,2,1} & \tuple{0,2,2} & \tuple{0,2,3} & \dots\\
\hline
\vdots & \vdots & \vdots & \vdots & \vdots & \ddots \\
\end{array}
\]
अशा रीतीने कँटर यांच्या नागमोडी पद्धतीसारखी पद्धत वापरून $\Nat^3$
चे प्रगणन मिळते. ही प्रक्रिया पुढे चालू ठेवून, प्रत्येक नैसर्गिक संख्या
$n$ साठी $\Nat^n$ चे प्रगणन मिळवता येते. म्हणून पुढील विधान मिळते:
\end{explain}

\begin{prop}
प्रत्येक $n \in \Nat$ साठी $\Nat^n$ हा गणनीय आहे.
\end{prop}

\begin{prob}\label{sfr:siz:zigzag:prob:posint-n}
प्रत्येक $n \in \Nat$ साठी $(\PosInt)^n$ हा गणनीय आहे, हे दाखवा.
\end{prob}

\begin{prob}\label{sfr:siz:zigzag:prob:posint-star}
  $(\PosInt)^*$ हा गणनीय आहे, हे दाखवा. तुम्ही \ref{sfr:siz:zigzag:prob:posint-n} गृहीत धरू शकता.
\end{prob}








\section{जोडीकरण फलने आणि संकेतांक}\label{sfr:siz:pai:sec}

\begin{explain}
कँटर यांच्या नागमोडी पद्धतीमुळे $\Nat^n$ ची गणनीयता चित्रातून
स्पष्ट दिसते. पण $\Nat^2$ दाखवणाऱ्या आपल्या सरणीवर लक्ष केंद्रित करू.
सरणीतील नागमोडी रेषेचा मागोवा घेऊन स्थाने मोजल्यास,
$\tuple{1,2}$ ही जोडी~$7$ या संख्येशी जोडली आहे हे तपासता येते.
पण ही संख्या अधिक थेट मोजता आली, तर सोयीचे होईल. म्हणजे, नागमोडी
प्रगणनाचे \emph{व्युत्क्रम} $g\colon \Nat^2 \to \Nat$ हाताशी असावे,
आणि त्यासाठी
\[g(\tuple{0,0}) = 0, \;
g(\tuple{0,1}) = 1, \;
g(\tuple{1,0}) = 2, \; \dots,
g(\tuple{1,2}) = 7, \; \dots
\]
असे असावे. त्यामुळे $\tuple{n, m}$ ही जोडी आपल्या प्रगणनात नेमकी
कोणत्या स्थानी येईल ते मोजता येईल.

खरे तर दोन निरीक्षणांच्या साहाय्याने $g$ ची थेट व्याख्या करता येते.
पहिले निरीक्षण: $n$ व्या ओळीतील आणि $m$ व्या स्तंभातील मूल्य~$v$
असेल, तर $(n+1)$ व्या ओळीतील आणि $(m-1)$ व्या स्तंभातील मूल्य
$v + 1$ असते. दुसरे निरीक्षण: आपल्या प्रगणनाची पहिली ओळ $0$, $1$,
$3$, $6$, इत्यादींपासून सुरू होणाऱ्या त्रिकोणी संख्यांची आहे.
$k$ वी त्रिकोणी संख्या म्हणजे $< k$ ही अट पूर्ण करणाऱ्या नैसर्गिक संख्यांची
बेरीज; ती $k(k+1)/2$ अशी मोजता येते. ही दोन निरीक्षणे एकत्र करून
पुढील फलन विचारात घ्या:
\[  g(n,m) = \frac{(n+m+1)(n+m)}{2} + n
\]
$g(\tuple{n, m})$ ऐवजी आपण बहुधा फक्त $g(n, m)$ असे लिहितो, कारण
ते पाहायला सोपे असते. या सूत्रानुसार आधी $(n+m)^\text{वी}$ त्रिकोणी
संख्या काढायची आणि नंतर तिच्यात $n$ मिळवायचा. त्यामुळे सरणी नेमकी
आपल्याला हवी तशी भरते. विशेषतः, $\tuple{1, 2}$ ही जोडी
$\frac{4 \times 3}{2} + 1 = 7$ कडे पाठवली जाते.

हे फलन $g$ म्हणजे जोड्यांच्या संचाच्या एका प्रगणनाचे \emph{व्युत्क्रम}
आहे. अशा फलनांना \emph{जोडीकरण फलने} म्हणतात.
\end{explain}

\begin{defn}[जोडीकरण फलन]
  $f\colon A \times B \to \Nat$ यातील $f$ हे एकास-एक असेल, तर ते
  \emph{अंकगणिती जोडीकरण फलन} असते. $f$ हे $A \times B$ चे
  \emph{सांकेतीकरण करते} आणि $f(x,y)$ हा $\tuple{x,y}$ चा
  \emph{संकेतांक} असतो असेही आपण म्हणतो.
\end{defn}

\begin{explain}
जोडीकरण फलनांनी, उदाहरणार्थ, नैसर्गिक संख्यांच्या जोड्यांचे
सांकेतीकरण करता येते; म्हणजेच घटकांची प्रत्येक \emph{जोडी} एका
\emph{एकाच} संख्येने दर्शवता येते. जोडीकरण फलनाचे व्युत्क्रम वापरून
संख्येचे \emph{विसांकेतीकरण} करता येते, म्हणजे ती संख्या कोणती जोडी
दर्शवते हे शोधता येते.
\end{explain}

\begin{prob}
सर्व ऋणेतर परिमेय संख्यांच्या संचाचे एक प्रगणन द्या.
\end{prob}

\begin{prob}
$\Rat$ हा गणनीय आहे, हे दाखवा. कोणतीही परिमेय संख्या
$z/m$ या अपूर्णांकाच्या रूपात लिहिता येते, जेथे $z \in \Int$ आणि
$m \in \Nat^+$ असतात, हे आठवा.
\end{prob}

\begin{prob}
$\Bin^*$ चे एक प्रगणन परिभाषित करा.
\end{prob}

\begin{prob}
प्रास्ताविक तर्कशास्त्राच्या अभ्यासक्रमातून आठवा की प्रत्येक संभाव्य
सत्यताकोष्टक एक सत्यता-फलन व्यक्त करतो. दुसऱ्या शब्दांत, सत्यता-फलने
म्हणजे कोणत्या तरी~$k$ साठी $\Bin^k \to \Bin$ अशी सर्व फलने आहेत.
सर्व सत्यता-फलनांचा संच गणनीय आहे, हे सिद्ध करा.
\end{prob}

\begin{prob}
कोणत्याही अनंत गणनीय संचाच्या सर्व सांत उपसंचांचा संच
गणनीय आहे, हे दाखवा.
\end{prob}

\begin{prob}
$\Nat$ च्या एखाद्या सांत उपसंचाचा पूरक असलेल्या $\Nat$ च्या उपसंचाला
\emph{सांत-पूरक} म्हणतात. म्हणजे, $A \subseteq \Nat$ हा सांत-पूरक
असतो तेव्हा आणि केवळ तेव्हाच, जेव्हा $\Nat\setminus A$ हा सांत असतो.
\begin{quote}\small\textbf{स्रोतदुरुस्ती.} OLSIZ-002: गोठवलेल्या इंग्रजी वाक्यात “complement of a finite set Nat” अशी संबंधदर्शक शब्दांशिवाय अपूर्ण रचना आहे. लगेचचे सूत्र अभिप्रेत अर्थ ठरवते: Nat मधील एखाद्या सांत उपसंचाचा Nat-सापेक्ष पूरक. मराठी मजकुरात हा संबंध स्पष्ट केला आहे.\end{quote}
$\Nat$ चे सांत आणि सांत-पूरक उपसंच हेच ज्याचे घटक आहेत,
असा संच $I$ घ्या. $I$ हा गणनीय आहे, हे दाखवा.
\end{prob}

\begin{prob}
गणनीय इतक्या गणनीय संचांचा संयोग गणनीय असतो,
हे दाखवा. म्हणजे, $A_1$, $A_2$, \dots{} हे संच असतील आणि प्रत्येक
$A_i$ हा गणनीय असेल, तर त्या सर्वांचा संयोग
$\bigcup_{i=1}^\infty A_i$ हाही गणनीय असतो. [टीप: हे कठीण आहे!]
\end{prob}

\begin{prob}
$f \colon A \times B \to \Nat$ हे कोणतेही जोडीकरण फलन असू द्या.
हे $f$ आच्छादकही असेल, तर त्याचे व्युत्क्रम $A \times B$ चे प्रगणन
आहे, हे दाखवा. सर्वसाधारण बाबतीत व्युत्क्रम केवळ फलनाच्या व्याप्तीवर
परिभाषित असते; यावरून जोड्यांचा हा संच गणनीय आहे, हे सिद्ध करा.
\begin{quote}\small\textbf{स्रोतदुरुस्ती.} MRSIZ-002: गोठवलेल्या इंग्रजी सरावात कोणत्याही एकास-एक जोडीकरण फलनाचे व्युत्क्रम थेट प्रगणन असल्याचा दावा आहे. फलन आच्छादक नसेल, तर त्याचे व्युत्क्रम सर्व नैसर्गिक संख्यांवर परिभाषित नसते. मराठी सरावात आच्छादक बाब आणि सामान्य अंशतः-व्युत्क्रम बाब वेगळ्या केल्या आहेत.\end{quote}
\end{prob}

\begin{prob}
$\Nat^3$ चे सांकेतीकरण करणारे फलन स्पष्टपणे द्या.
\end{prob}









\section{एक पर्यायी जोडीकरण फलन}\label{sfr:siz:pai-alt:sec}

\begin{explain}
$\Nat^2$ ची अशी इतर प्रगणनेही आहेत की त्यांची व्युत्क्रमे शोधणे अधिक
सोपे होते. त्यांपैकी एक पुढे दिले आहे. प्रगणन सरणीत दाखवण्याऐवजी,
सुरुवातीला रिकाम्या असलेल्या जागांशी जोडलेल्या धन पूर्णांकांच्या
यादीपासून सुरुवात करा. या जागा पुढीलप्रमाणे एकामागून एक
$\tuple{n,m}$ जोड्यांनी भरत आहोत, अशी कल्पना करा. पहिल्या स्थानी~$0$
असलेल्या जोड्यांपासून (म्हणजे $\tuple{0,m}$ जोड्यांपासून) सुरुवात करा.
पहिली जोडी (म्हणजे $\tuple{0,0}$) पहिल्या रिकाम्या जागेत ठेवा; मग एक
रिकामी जागा सोडून दुसरी जोडी पुढील रिकाम्या जागेत ठेवा आणि पुन्हा एक
जागा सोडा; ही प्रक्रिया अशीच पुढे चालू ठेवा. गोठवलेल्या इंग्रजी
मजकुरात दुसरी जोडी $\tuple{0,2}$ अशी दिली आहे; परंतु खालील सारणीनुसार
येथे $\langle 0,1\rangle$ अभिप्रेत आहे. आता आपल्या प्रगणनाची (अपूर्ण) सुरुवात अशी
दिसते:
\begin{quote}\small\textbf{स्रोतदुरुस्ती.} MRSIZ-003: गोठवलेल्या इंग्रजी मजकुरात एक रिकामी जागा सोडल्यानंतरची दुसरी जोडी \ensuremath{\langle}0,2\ensuremath{\rangle} अशी दिली आहे, पण लगेचची सारणी आणि वर्णनाची क्रमवार रचना ती \ensuremath{\langle}0,1\ensuremath{\rangle} असावी हे दाखवतात. सूत्र-संरचना जशीच्या तशी ठेवून मराठी मजकुरात अपेक्षित जोडी स्पष्ट केली आहे.\end{quote}
\[\begin{array}{@{}c c c c c c c c c c c@{}}
\mathbf 1 & \mathbf 2 & \mathbf 3 & \mathbf 4 & \mathbf 5 & \mathbf 6 & \mathbf 7 & \mathbf 8 & \mathbf 9 & \mathbf{10} & \dots \\ \\
\tuple{0,0} &  & \tuple{0,1} &  & \tuple{0,2} &  & \tuple{0,3} & & \tuple{0,4} &  & \dots \\
\end{array}
\]
अजून रिकाम्या राहिलेल्या जागांसाठी $\tuple{1,m}$ जोड्यांबाबत हीच
प्रक्रिया पुन्हा करा; पुन्हा प्रत्येक दुसरी रिकामी जागा वगळा:
\[\begin{array}{@{}c c c c c c c c c c c@{}}
\mathbf 1 & \mathbf 2 & \mathbf 3 & \mathbf 4 & \mathbf 5 & \mathbf 6 & \mathbf 7 & \mathbf 8 & \mathbf 9 & \mathbf{10} & \dots \\ \\
\tuple{0,0} & \tuple{1,0} & \tuple{0,1} &  & \tuple{0,2} & \tuple{1,1} &
\tuple{0,3} & & \tuple{0,4} &  \tuple{1,2} & \dots \\
\end{array}
\]
यानंतर $\tuple{2,m}$ जोड्या आणि मग $\tuple{3,m}$ जोड्या त्याच पद्धतीने
भरा; त्यापुढेही पहिला घटक अनुक्रमे 4, 5, इत्यादी करत पुढे जायचे आहे.
आपल्या पूर्ण प्रगणनाची सुरुवात म्हणून अशी दिसते:
\begin{quote}\small\textbf{स्रोतदुरुस्ती.} OLSIZ-003: गोठवलेल्या इंग्रजी मजकुरात “pairs \ensuremath{\langle}2,m\ensuremath{\rangle}, \ensuremath{\langle}2,m\ensuremath{\rangle}, etc.” अशी पुनरुक्ती आहे. पूर्ण सारणीतील \ensuremath{\langle}3,0\ensuremath{\rangle} आणि पुढील सर्व ओळींवरून दुसऱ्या ठिकाणी \ensuremath{\langle}3,m\ensuremath{\rangle} अभिप्रेत असल्याचे निश्चित होते; मराठी मजकुरात दुसरी जोडी त्यानुसार दुरुस्त केली आहे.\end{quote}
\[\begin{array}{@{}cc c c c c c c c c c@{}}
\mathbf 1 & \mathbf 2 & \mathbf 3 & \mathbf 4 & \mathbf 5 & \mathbf 6 & \mathbf 7 & \mathbf 8 & \mathbf 9 & \mathbf{10} & \dots \\ \\
\tuple{0,0} & \tuple{1,0} & \tuple{0,1} & \tuple{2,0}  & \tuple{0,2} &
\tuple{1,1} & \tuple{0,3} & \tuple{3,0}  & \tuple{0,4} &  \tuple{1,2} & \dots \\
\end{array}
\]
वरील सरणीतील चौकटींना या प्रगणनानुसार क्रमांक दिले, तर आपल्याला
नीट नागमोडी रेषा दिसणार नाही; त्याऐवजी पुढील मांडणी दिसेल:
\[\begin{array}{ c | c | c | c | c | c | c | c }
& \mathbf 0 & \mathbf 1 & \mathbf 2 & \mathbf 3 & \mathbf 4 & \mathbf 5 & \dots \\
\hline
\mathbf 0 & 1 & 3 & 5 & 7 & 9 & 11 & \dots \\
\hline
\mathbf 1 & 2 & 6 & 10 & 14 & 18 & \dots & \dots \\
\hline
\mathbf 2 & 4 & 12 & 20 & 28 & \dots & \dots & \dots \\
\hline
\mathbf 3 & 8 & 24 & 40 & \dots & \dots & \dots & \dots \\
\hline
\mathbf 4 & 16 & 48 & \dots & \dots & \dots & \dots & \dots \\
\hline
\mathbf 5 & 32 & \dots & \dots & \dots & \dots & \dots & \dots \\
\hline
\vdots & \vdots & \vdots & \vdots & \vdots & \vdots & \vdots & \ddots\\
\end{array}
\]
आपल्या प्रगणनात ओळ~$0$ मधील जोड्या विषम क्रमांकांच्या स्थानी आहेत,
हे दिसते; म्हणजे $\tuple{0,m}$ ही जोडी $2m+1$ या स्थानी आहे. दुसऱ्या
ओळीतील $\tuple{1,m}$ जोड्या विषम संख्येच्या दुप्पट क्रमांक असलेल्या
स्थानी आहेत; विशेषतः त्या $2 \cdot (2m+1)$ या स्थानी आहेत. तिसऱ्या
ओळीतील $\tuple{2,m}$ जोड्या विषम संख्येच्या चौपट क्रमांक असलेल्या
स्थानी, म्हणजे $4 \cdot (2m+1)$ या स्थानी आहेत; आणि ही रचना अशीच
पुढे चालते. प्रत्येक ओळीत $(2m+1)$ च्या गुणकांचे मूल्य $1$, $2$,
$4$, $8$, \dots{} असे आहे; हे नेमके~$2$ चे घात आहेत:
$1= 2^0$, $2 = 2^1$, $4 = 2^2$, $8 = 2^3$, \dots\@ विचाराधीन
जोडीचा पहिला घटक हाच नेहमी संबंधित घातांक असतो. म्हणून
$\tuple{n,m}$ या जोडीसाठी गुणक $2^n$ आहे. यातून आपल्याला
$2^n \cdot (2m+1)$ हे सर्वसाधारण सूत्र मिळते. पण या फलनाने जोड्या
\emph{धन} पूर्णांकांकडे पाठवल्या जातात; म्हणजे $\tuple{0,0}$ चे
स्थान~$1$ आहे. स्थानांची सुरुवात~$0$ पासून करायची असेल, तर उत्तरातून
$1$ वजा करावा लागतो. त्यामुळे पुढील फलन मिळते:
\end{explain}

\begin{ex}
$h\colon \Nat^2 \to \Nat$ हे पुढीलप्रमाणे दिलेले फलन घ्या:
\[h(n,m) = 2^n (2m+1) - 1
\]
हे नैसर्गिक संख्यांच्या जोड्यांचा संच~$\Nat^2$ यासाठीचे जोडीकरण फलन
आहे.
\end{ex}

\begin{explain}
त्यानुसार, $\Nat^2$ च्या आपल्या दुसऱ्या प्रगणनात $\tuple{0,0}$ या
जोडीचा संकेतांक $h(0,0) = 2^0(2\cdot 0+1) - 1 = 0$ आहे;
$\tuple{1,2}$ हिचा संकेतांक $2^{1} \cdot (2 \cdot 2 + 1) - 1 = 2
\cdot 5 - 1 = 9$ आहे; आणि $\tuple{2,6}$ हिचा संकेतांक $2^{2} \cdot (2
\cdot 6 + 1) - 1 = 51$ आहे.
\end{explain}

कधी कधी नैसर्गिक संख्यांच्या जोड्यांचे~$\Nat^2$ सांकेतीकरण करताना ते
आच्छादक असण्याची आवश्यकता नसते. अशा सांकेतीकरणांची व्युत्क्रमे केवळ
अंशतः फलने असतात.

\begin{ex}
$j\colon \Nat^2 \to \Nat^+$ हे पुढीलप्रमाणे दिलेले फलन घ्या:
\[j(n,m) = 2^n3^m
\]
हे $\Nat^2 \to \Nat$ असे एकास-एक फलन आहे.
\end{ex}









\section{अगणनीय संच}\label{sfr:siz:nen:sec}

\begin{editorial}
  या विभागात \ref{sfr:siz:enm:sec} मधील व्याख्या वापरून $\Bin^\omega$ आणि
  $\Pow{\PosInt}$ यांची अगणनीयता सिद्ध केली आहे. ही मांडणी
  \ref{sfr:siz:enm-alt:sec} मधील आवृत्तीपेक्षा थोडी अधिक प्राथमिक आणि
  थोडी अधिक तपशीलवार ठेवली आहे.
\end{editorial}

धन पूर्णांकांचा संच~$\PosInt$ यासारखे काही संच अनंत असतात. आतापर्यंत
पाहिलेल्या अनंत संचांची सर्व उदाहरणे गणनीय होती. पण हा गुणधर्म
नसलेले अनंत संचही असतात. अशा संचांना \emph{अगणनीय} म्हणतात.

सुरुवातीलाच, अगणनीय संच असतात ही गोष्ट कदाचित आश्चर्यकारक
वाटेल. कोणत्याही गणनीय संच~$A$ साठी $f \colon \PosInt \to A$
असे आच्छादक फलन असते. संच अगणनीय असेल, तर असे फलन
असू शकत नाही. म्हणजे, $\PosInt$ मधील अनंत घटक ना~$A$ कडे
पाठवणारे कोणतेही फलन~$A$ मधील सर्व घटक व्यापू शकत नाही. या अर्थाने,
अनंत असलेल्या धन पूर्णांकांपेक्षा~$A$ मध्ये ``अधिक'' घटक असतात.

एखादा संच अगणनीय आहे हे कसे सिद्ध करायचे? तसे कोणतेही
आच्छादक फलन अस्तित्वात असू शकत नाही, हे दाखवावे लागते. समतुल्यपणे,
$A$ मधील घटकांचे एकदिश अनंत यादीत प्रगणन करता येत नाही, हे दाखवावे
लागते. हे करण्याचा सर्वोत्तम मार्ग म्हणजे~$A$ मधील घटक च्या
प्रत्येक यादीतून किमान एक घटक सुटतो, किंवा $f\colon \PosInt \to A$
असे कोणतेही फलन आच्छादक असू शकत नाही, हे दाखवणे. यासाठी कँटर
यांची \emph{कर्णरेषीय पद्धत} वापरता येते. $A$ मधील घटक ची
एखादी यादी, म्हणजे $x_1$, $x_2$, \dots, दिली असता, आपण~$A$ मधील असा
दुसरा घटक रचतो की त्याच्या रचनेमुळे तो त्या यादीत असणे अशक्य ठरते.

आपले पहिले उदाहरण म्हणजे $0$ आणि $1$ यांच्या सर्व अनंत व मध्ये खंड
नसलेल्या अनुक्रमांचा संच~$\Bin^\omega$.

\begin{thm}
\label{sfr:siz:nen:thm:nonenum-bin-omega}
$\Bin^\omega$ हा अगणनीय आहे.
\end{thm}

\begin{proof}
व्याघाताने सिद्ध करण्यासाठी असे समजा की $\Bin^\omega$ हा गणनीय आहे; म्हणजे,
$\Bin^\omega$ मधील सर्व घटक ची $s_{1}$, $s_{2}$, $s_{3}$,
$s_{4}$, \dots{} अशी यादी आहे, असे समजा. यांपैकी प्रत्येक $s_i$ हा
स्वतः $0$ आणि~$1$ यांचा अनंत अनुक्रम आहे. या यादीतील $i$ व्या अनुक्रमाच्या
$j$ व्या घटकाला $s_i(j)$ म्हणू. मग $i$ वा अनुक्रम~$s_i$ असा आहे:
\[s_i(1), s_i(2), s_i(3), \dots
\]

ही यादी आणि तिच्यातील प्रत्येक अनुक्रम~$s_i$ चे घटक आपण पुढील सरणीत
मांडू शकतो:
\[\begin{array}{c|c|c|c|c|c}
& 1 & 2 & 3 & 4 & \dots \\\hline
1 & \mathbf{s_{1}(1)} & s_{1}(2) & s_{1}(3) & s_1(4) & \dots \\\hline
2 & s_{2}(1)& \mathbf{s_{2}(2)} & s_2(3) & s_2(4) & \dots \\\hline
3 & s_{3}(1)& s_{3}(2) & \mathbf{s_3(3)} & s_3(4) & \dots \\\hline
4 & s_{4}(1)& s_{4}(2) & s_4(3) & \mathbf{s_4(4)} & \dots \\\hline
\vdots & \vdots & \vdots & \vdots & \vdots & \mathbf{\ddots}
\end{array}
\]
डाव्या बाजूकडील खुणा यादीतील अनुक्रमांना $s_1$, $s_2$, \dots{} असे
क्रमांक देतात; वरच्या बाजूकडील संख्या प्रत्येक अनुक्रमातील घटक
ना क्रमांक देतात. उदाहरणार्थ, $s_{1}(1)$ हे अनुक्रम~$s_{1}$ मधील
पहिला घटक असलेल्या संख्येचे नाव आहे; ती संख्या $0$ किंवा~$1$
असू शकते. पुढेही याचप्रमाणे अर्थ घ्यायचा.

आता आपण $0$ आणि $1$ यांचा असा अनंत अनुक्रम~$\overline{s}$ रचू, जो
या यादीत असणे अशक्य आहे. $\overline{s}$ ची व्याख्या $s_1$, $s_2$,
\dots{} या यादीवर अवलंबून असेल. $0$ आणि $1$ यांच्या अनंत अनुक्रमांची
कोणतीही अनंत यादी असा अनंत अनुक्रम~$\overline{s}$ निर्माण करते, जो
त्या यादीत नसेल याची खात्री असते.

$\overline{s}$ ची व्याख्या करण्यासाठी आपण त्याचे सर्व घटक
ठरवतो; म्हणजे, प्रत्येक $n \in \PosInt$ साठी $\overline{s}(n)$
ठरवतो. वरील सरणीच्या कर्णरेषेवरून खाली वाचून (म्हणूनच ``कर्णरेषीय
पद्धत'' हे नाव) प्रत्येक $1$ चे $0$ मध्ये आणि प्रत्येक $0$ चे~$1$
मध्ये रूपांतर करून आपण हे करतो. अधिक अमूर्तपणे, कर्णरेषेवरील $n$ वा
घटक, म्हणजे $s_n(n)$, हा $1$ आहे की $0$ यानुसार
$\overline{s}(n)$ हे $0$ किंवा $1$ असे परिभाषित करतो:
\[\overline{s}(n) =
\begin{cases}
1 & \text{जर $s_{n}(n) = 0$ असेल}\\
0 & \text{जर $s_{n}(n) = 1$ असेल}.
\end{cases}
\]
प्रकरणांनुसार व्याख्येपेक्षा सूत्रे अधिक आवडत असतील, तर
$\overline{s}(n) = 1 - s_n(n)$ अशीही व्याख्या करता येईल.

$\overline{s}$ हा स्पष्टपणे $0$ आणि $1$ यांचा अनंत अनुक्रम आहे, कारण
तो आपल्या सरणीच्या कर्णरेषेवर दिसणाऱ्या $0$ आणि $1$ यांच्या अनुक्रमाचा
पूरक अनुक्रम आहे. म्हणून $\overline{s}$ हा $\Bin^\omega$ चा
घटक आहे. पण तो $s_1$, $s_2$, \dots{} या यादीत असू शकत नाही.
असे का?

तो यादीतील पहिला अनुक्रम~$s_1$ असू शकत नाही, कारण त्याचा पहिला
घटक हा $s_1$ च्या पहिल्या घटकापेक्षा वेगळा आहे. $s_1(1)$
काहीही असो, $\overline{s}(1)$ त्याच्या विरुद्ध असेल अशी आपण व्याख्या
केली. तो यादीतील दुसरा अनुक्रमही असू शकत नाही, कारण $\overline{s}$
चा दुसरा घटक $s_2$ च्या दुसऱ्या घटकापेक्षा वेगळा आहे: $s_2(2)$ हा
$0$ असेल, तर $\overline{s}(2)$ हा $1$ आहे, आणि उलट. पुढेही असेच.

अधिक नेमकेपणाने सांगायचे तर, $\overline{s}$ यादीत असता, तर कोणत्या
तरी $k$ साठी $\overline{s} = s_{k}$ असते. दोन अनुक्रम प्रत्येक स्थानी
जुळतात तेव्हा आणि केवळ तेव्हाच ते एकरूप असतात; म्हणजे, कोणत्याही~$n$
साठी $\overline{s}(n) = s_{k}(n)$ असते. म्हणून, विशेष प्रकरण म्हणून
$n = k$ घेतल्यास $\overline{s}(k) = s_{k}(k)$ असणे आवश्यक ठरते.
$s_k(k)$ हे $0$ किंवा~$1$ आहे. ते $0$ असेल, तर $\overline{s}(k)$ हे~$1$
असले पाहिजे---कारण आपण $\overline{s}$ ची तशीच व्याख्या केली आहे. पण
$s_k(k) = 1$ असेल, तर पुन्हा $\overline{s}$ च्या व्याख्येमुळे
$\overline{s}(k) = 0$ होते. दोन्ही बाबतीत
$\overline{s}(k) \neq s_{k}(k)$.

सुरुवातीला आपण $\Bin^\omega$ मधील घटक ची $s_1$, $s_2$,
\dots{} अशी यादी आहे असे गृहीत धरले. या यादीपासून आपण असा
अनुक्रम~$\overline{s}$ रचला की तो यादीत असू शकत नाही, हे सिद्ध केले.
पण सर्व $s_i$ हे $0$ आणि $1$ यांचे अनुक्रम असतील, तर तो नक्कीच $0$
आणि $1$ यांचा अनुक्रम असतो; म्हणजे, $\overline{s} \in
\Bin^\omega$. यावरून विशेषतः $\Bin^\omega$ मधील \emph{सर्व}
घटक ची यादी असू शकत नाही. कारण अशी कोणतीही यादी दिली, तरी
तिच्यात नसेल असा अनुक्रम~$\overline{s}$ आपण रचू शकतो. म्हणून
$\Bin^\omega$ मधील सर्व अनुक्रमांची यादी आहे हे गृहीतक व्याघात निर्माण
करते.
\end{proof}

\begin{explain}
या सिद्धता-पद्धतीला ``कर्णीकरण'' म्हणतात, कारण तिच्यात सरणीच्या
कर्णरेषेचा वापर करून~$\overline{s}$ ची व्याख्या केली जाते. कर्णीकरणात
सरणी असलीच पाहिजे असे नाही: प्रत्यक्ष सरणी किंवा कर्णरेषा नसतानाही
हीच कल्पना वापरून संच गणनीय नाहीत, हे दाखवता येते.
\end{explain}

\begin{thm}
\label{sfr:siz:nen:thm:nonenum-pownat}
$\Pow{\PosInt}$ हा गणनीय नाही.
\end{thm}

\begin{proof}
आपण त्याच पद्धतीने पुढे जाऊ: $\PosInt$ च्या उपसंचांची प्रत्येक यादी
दिली असता त्या यादीत नसलेला $\PosInt$ चा एक उपसंच असतो, हे दाखवू.
$\PosInt$ च्या उपसंचांची पुढील यादी दिली आहे असे समजा:
\[Z_{1}, Z_{2}, Z_{3}, \dots
\]
आता आपण $\overline{Z}$ हा असा संच परिभाषित करतो की प्रत्येक
$n \in \PosInt$ साठी $n \in \overline{Z}$ तेव्हा आणि केवळ तेव्हाच,
जेव्हा $n \notin Z_{n}$:
\[\overline{Z} = \Setabs{n \in \PosInt}{n \notin Z_n}
\]
$\overline{Z}$ हा स्पष्टपणे धन पूर्णांकांचा संच आहे, कारण गृहीतकानुसार
प्रत्येक~$Z_n$ तसा आहे; म्हणून $\overline{Z} \in
\Pow{\PosInt}$. पण $\overline{Z}$ यादीत असू शकत नाही. हे दाखवण्यासाठी
प्रत्येक $k \in \PosInt$ साठी $\overline{Z} \neq Z_k$, हे सिद्ध करू.

आता $k \in \PosInt$ स्वैर असू द्या. प्रत्येक $n \in \PosInt$ साठी
$n \in \overline{Z}$ तेव्हा आणि केवळ तेव्हाच, जेव्हा $n \notin Z_n$,
अशी $\overline{Z}$ ची व्याख्या केली आहे. विशेषतः $n=k$ घेतल्यास,
$k \in \overline{Z}$ तेव्हा आणि केवळ तेव्हाच, जेव्हा $k \notin Z_k$.
पण यावरून $\overline{Z} \neq Z_k$ हे दिसते, कारण $k$ हा एकाचा
घटक आहे आणि दुसऱ्याचा नाही; म्हणून $\overline{Z}$ आणि $Z_k$
यांचे घटक वेगवेगळे आहेत. $k$ स्वैर असल्याने $\overline{Z}$ हा
$Z_1$, $Z_2$, \dots{} या यादीत नाही.
\end{proof}

\begin{explain}
मागील सिद्धतेत कर्णरेषेचा उल्लेख झाला नाही; पण पुढील चित्र मनात आणले,
तर तिच्यात कर्णरेषा आहे असे समजता येते. $Z_1$, $Z_2$, \dots{} हे संच
एका सरणीत लिहिले आहेत अशी कल्पना करा; प्रत्येक घटक~$j \in Z_i$
हा $j$ व्या स्तंभात लिहिला आहे. समजा, त्या यादीतील पहिले चार संच
$\{1,2,3,\dots\}$, $\{2, 4, 6, \dots\}$, $\{1,2,5\}$ आणि
$\{3,4,5,\dots\}$ आहेत. मग सरणीची सुरुवात अशी होईल:
\[\begin{array}{r@{}rrrrrrr}
  Z_1 = \{ & \mathbf{1}, & 2, & 3, & 4, & 5, & 6, & \dots\}\\
  Z_2 = \{ &  & \mathbf{2}, &  & 4, &  & 6, & \dots\}\\
  Z_3 = \{ & 1, & 2, &  &  & 5\phantom{,} &  & \}\\
  Z_4 = \{ &  &  & 3, & \mathbf{4}, & 5, & 6, & \dots\}\\
  \vdots & & & & & \ddots
\end{array}
\]
यानंतर कर्णरेषेवरून खाली जाताना कर्णरेषेवर दिसणाऱ्या संख्या वगळून,
आणि $j$ व्या ओळीच्या त्याच क्रमांकाच्या स्तंभातील जागा रिकामी असेल अशा~$j$
चा समावेश करून, $\overline{Z}$ हा संच मिळतो. वरील उदाहरणात
आपण $1$ आणि $2$ वगळू, $3$ चा समावेश करू, $4$ वगळू, इत्यादी.
\end{explain}

\begin{prob}
कर्णरेषीय युक्तिवादाने $\Pow{\Nat}$ हा अगणनीय आहे, हे दाखवा.
\end{prob}

\begin{prob}\label{sfr:siz:nen:prob:f-posint}
$f \colon \PosInt \to \PosInt$ अशा फलनांचा संच अगणनीय
आहे, हे स्पष्ट कर्णरेषीय युक्तिवादाने दाखवा. म्हणजे, $f_1$, $f_2$,
\dots{} ही फलनांची यादी असेल आणि प्रत्येक $f_i\colon
\PosInt \to \PosInt$ असेल, तर या यादीत नसलेले कोणते तरी
$\overline{f}\colon \PosInt \to \PosInt$ असते, हे दाखवा.
\end{prob}









\section{न्यूनीकरण}\label{sfr:siz:red:sec}

\begin{editorial}
  या विभागात \ref{sfr:siz:nen:sec} मधील निष्कर्षांशी जुळणारी अगणनीयता
  न्यूनीकरणाने सिद्ध केली आहे. \ref{sfr:siz:nen-alt:sec} मधील निष्कर्षांशी
  जुळणारी, किंचित अधिक संक्षिप्त पर्यायी आवृत्ती \ref{sfr:siz:red-alt:sec}
  मध्ये दिली आहे.
\end{editorial}

कर्णीकरणाच्या युक्तिवादाने $\Pow{\PosInt}$ हा अगणनीय आहे,
हे आपण दाखवले. सर्व अनंत $0$--$1$ अनुक्रमांचा संच~$\Bin^\omega$ हा
अगणनीय आहे, याची सिद्धता आपल्याकडे आधीच होती.
$\Pow{\PosInt}$ हा अगणनीय आहे हे सिद्ध करण्याचा आणखी एक मार्ग
असा: \emph{$\Pow{\PosInt}$ हा गणनीय असेल, तर $\Bin^\omega$
  हाही गणनीय असतो} हे दाखवा. $\Bin^\omega$ हा गणनीय
नाही, हे माहीत असल्याने $\Pow{\PosInt}$ हाही तसा असू शकत नाही. यालाच
एक समस्या दुसऱ्या समस्येकडे \emph{न्यूनीत करणे} म्हणतात---या बाबतीत
$\Bin^\omega$ चे प्रगणन करण्याची समस्या आपण $\Pow{\PosInt}$ चे प्रगणन
करण्याच्या समस्येकडे न्यूनीत करतो. नंतरच्या समस्येचे उत्तर---म्हणजे
$\Pow{\PosInt}$ चे प्रगणन---पहिल्या समस्येचे उत्तर---म्हणजे
$\Bin^\omega$ चे प्रगणन---देईल.

संच~$B$ च्या प्रगणनाची समस्या संच~$A$ च्या प्रगणनाच्या समस्येकडे कशी
न्यूनीत करायची? $A$ चे प्रगणन $B$ च्या प्रगणनात रूपांतरित करण्याची
पद्धत आपण देतो. त्यासाठीचा सर्वांत सोपा मार्ग म्हणजे
$f\colon A \to B$ असे आच्छादक फलन परिभाषित करणे. $x_1$, $x_2$,
\dots{} हे~$A$ चे प्रगणन असेल, तर $f(x_1)$, $f(x_2)$, \dots{} हे~$B$
चे प्रगणन करील. आपल्या बाबतीत आपण
$f\colon \Pow{\PosInt} \to \Bin^\omega$ असे आच्छादक फलन शोधत आहोत.

\begin{prob}
$g\colon B \to A$ असे एकास-एक फलन असेल आणि $B$ हा
अगणनीय असेल, तर $A$ हाही अगणनीय असतो, हे दाखवा.
$g$ वापरून~$A$ चे प्रगणन~$B$ च्या प्रगणनात कसे रूपांतरित करता येते,
हे दाखवून सिद्धता द्या.
\end{prob}

\begin{proof}[न्यूनीकरणाने {\ref{sfr:siz:nen:thm:nonenum-pownat}} ची सिद्धता]
$\Pow{\PosInt}$ हा गणनीय आहे, आणि म्हणून त्याचे $Z_{1}$,
$Z_{2}$, $Z_{3}$, \dots{} असे प्रगणन आहे, असे समजा.

$f \colon \Pow{\PosInt} \to \Bin^\omega$ हे फलन पुढीलप्रमाणे परिभाषित
करा: $f(Z)$ हा असा अनुक्रम~$s$ असू द्या की $n \in Z$ असेल तेव्हा
आणि केवळ तेव्हाच $s(n) = 1$, आणि अन्यथा $s(n) = 0$. याने
फलन स्पष्टपणे परिभाषित होते, कारण $Z \subseteq \PosInt$ असताना कोणताही
$n \in \PosInt$ हा $Z$ चा घटक असतो किंवा नसतो. उदाहरणार्थ, धन सम संख्यांचा संच
$2\PosInt = \{2, 4, 6, \dots\}$ हा $010101\dots$ या अनुक्रमाकडे,
रिक्त संच $0000\dots$ कडे आणि $\PosInt$ हा संच स्वतः $1111\dots$ कडे
पाठवला जातो.
\begin{quote}\small\textbf{स्रोतदुरुस्ती.} OLSIZ-004: गोठवलेल्या इंग्रजी व्याख्येत f(Z) ला s\_k असे म्हटले आहे, पण k चा Z शी संबंध किंवा निवड दिलेली नाही. अभिप्रेत अनुक्रम Z नेच ठरतो; मराठी व्याख्येत एकच बद्ध निर्गत नाव s आणि त्याचे घटक s(n) सातत्याने वापरले आहेत.\end{quote}

हे फलन आच्छादक सुद्धा आहे: $0$ आणि $1$ यांचा प्रत्येक अनुक्रम धन
पूर्णांकांच्या कोणत्या तरी संचाशी जुळतो; तो संच म्हणजे अनुक्रमात ज्या
स्थानी~$1$ आहे त्या स्थानांना अनुरूप पूर्णांकांचा संच. अधिक नेमकेपणाने,
$s \in \Bin^\omega$ असे समजा. पुढीलप्रमाणे $Z \subseteq \PosInt$
परिभाषित करा:
\[Z = \Setabs{n \in \PosInt}{s(n) = 1}
\]
मग $f$ च्या व्याख्येवरून तपासता येते की $f(Z) = s$.

आता पुढील यादी विचारात घ्या:
\[f(Z_1), f(Z_2), f(Z_3), \dots
\]
$f$ हे आच्छादक असल्याने $\Bin^\omega$ मधील प्रत्येक सदस्य हा
$f$ च्या कोणत्या तरी फलसाधकावरील मूल्य असला पाहिजे, आणि म्हणून तो
यादीत आला पाहिजे. त्यामुळे ही यादी $\Bin^\omega$ मधील सर्व घटकांचे
प्रगणन करते.

म्हणून $\Pow{\PosInt}$ हा गणनीय असता, तर $\Bin^\omega$ हाही
गणनीय असता. पण $\Bin^\omega$ हा अगणनीय आहे
(\ref{sfr:siz:nen:thm:nonenum-bin-omega}). त्यामुळे $\Pow{\PosInt}$ हा
अगणनीय आहे.
\end{proof}

\begin{explain}
न्यूनीकरण कोणत्या दिशेने केले जाते याबद्दल गोंधळ होणे सोपे आहे.
उदाहरणार्थ, $g \colon \Bin^\omega \to B$ असे आच्छादक फलन
$B$ हा अगणनीय आहे, हे \emph{सिद्ध करत नाही}. ($g
\colon \Bin^\omega \to \Bin$ हे $g(s) = s(1)$ ने परिभाषित केलेले,
$0$ आणि $1$ यांच्या अनुक्रमाला त्याच्या पहिल्या घटक कडे
पाठवणारे फलन विचारात घ्या. काही अनुक्रम $0$ ने आणि काही $1$ ने सुरू
होत असल्याने ते आच्छादक आहे. पण $\Bin$ हा सांत संच आहे.)
तसेच फलन~$f$ हे आच्छादक असलेच पाहिजे; अन्यथा युक्तिवाद लागू
होत नाही: $f(x_1)$, $f(x_2)$, \dots{} यात~$B$ मधील सर्व घटक
असतील याची खात्री उरत नाही. उदाहरणार्थ,
\[h(n) = \underbrace{000\dots0}_{\text{$n$ वेळा $0$}}111\dots
\]
याने $h\colon \PosInt \to
\Bin^\omega$ असे फलन परिभाषित होते: प्रत्येक मूल्याची सुरुवात n शून्यांनी
होते आणि त्यानंतर कायम एक येतात. या फलनाच्या व्याप्तीत अखेरीस कायम एक
न होणारे अनुक्रम येत नाहीत, म्हणून ते आच्छादक नाही; पण $\PosInt$ हा
गणनीय आहे.
\begin{quote}\small\textbf{स्रोतदुरुस्ती.} OLSIZ-005: गोठवलेल्या इंग्रजी प्रदर्शनात h(n) चे मूल्य फक्त n शून्यांची सांत चिन्हमाला आहे, जी अनंत अनुक्रमांच्या Bin-omega संचात येत नाही. मराठी प्रदर्शनात त्यानंतर अनंत एकांची शेपटी जोडून वैध, अनाच्छादक फलन दिले आहे.\end{quote}
\end{explain}

\begin{prob}
धन पूर्णांकांच्या जोड्यांच्या \emph{सर्व संचांचा} संच न्यूनीकरणाच्या
युक्तिवादाने अगणनीय आहे, हे दाखवा.
\end{prob}

\begin{prob}\label{sfr:siz:red:prob:nat-nat}
  $f\colon \Nat \to \Nat$ अशा सर्व फलनांचा संच~$X$ हा न्यूनीकरणाच्या
  युक्तिवादाने अगणनीय आहे, हे दाखवा. (सूचना: $X$ पासून
  $\Bin^\omega$ कडे जाणारे आच्छादक फलन द्या.)
\end{prob}

\begin{prob}
नैसर्गिक संख्यांच्या सर्व अनंत अनुक्रमांचा संच~$\Nat^\omega$ हा
न्यूनीकरणाच्या युक्तिवादाने अगणनीय आहे, हे दाखवा.
\end{prob}

\begin{prob}
$P$ हा धन पूर्णांकांच्या संचापासून $\{0\}$ या संचाकडे जाणाऱ्या फलनांचा
संच असू द्या; आणि $Q$ हा धन पूर्णांकांच्या संचापासून $\{0\}$ या
संचाकडे जाणाऱ्या \emph{अंशतः} फलनांचा संच असू द्या. $P$ हा
गणनीय आहे आणि $Q$ नाही, हे दाखवा. (सूचना: $\Bin^\omega$ चे
प्रगणन करण्याची समस्या~$Q$ चे प्रगणन करण्याच्या समस्येकडे न्यूनीत करा.)
\end{prob}

\begin{prob}
$S$ हा धन पूर्णांकांच्या संचापासून \{0,1\} या संचाकडे जाणाऱ्या सर्व
आच्छादक फलनांचा संच असू द्या; म्हणजे, $S$ मध्ये सर्व
आच्छादक~$f\colon \PosInt \to \Bin$ आहेत. $S$ हा
अगणनीय आहे, हे दाखवा.
\end{prob}

\begin{prob}
सर्व वास्तव संख्यांचा संच~$\Real$ हा अगणनीय आहे, हे दाखवा.
\end{prob}









\section{तुल्यबलता}\label{sfr:siz:equ:sec}

संचांच्या ``आकारमाना''ची एक अंतःप्रज्ञ कल्पना आपल्याकडे आहे आणि ती
सांत संचांसाठी व्यवस्थित चालते. पण अनंत संचांचे काय? \emph{कोणत्याही}
आकारमानाच्या दोन संचांच्या आकारमानांची औपचारिक रीतीने तुलना करायची
असेल, तर दोन संच एकाच आकारमानाचे कधी असतात याची व्याख्या प्रथम करणे
योग्य ठरते. फ्रेगे यांचे पुढील उदाहरण पाहा:
\begin{quote}
  उपाहारगृहातील कर्मचाऱ्याला मेजावर ताटांइतक्याच सुऱ्या मांडल्या आहेत
  याची खात्री करायची असेल, तर प्रत्येक ताटाच्या उजवीकडे एक सुरी ठेवून
  आणि मेजावरील प्रत्येक सुरी कोणत्या तरी ताटाच्या उजवीकडे येईल असे
  केल्यास त्याला ताटे किंवा सुऱ्या मोजण्याची गरज नसते. अशा प्रकारे
  ताटे आणि सुऱ्या एकास-एक संगतीने परस्परांशी जोडली जातात, आणि तीही
  त्याच अवकाशीय संबंधाने. (Frege, 1884, \S70)
\end{quote}
या उताऱ्यातील अंतर्दृष्टी पुढील औपचारिक व्याख्येत मांडता येते:

\begin{defn}\label{sfr:siz:equ:comparisondef}
  $A$ हा $B$ शी \emph{तुल्यबल} असतो, आणि ते $\cardeq{A}{B}$ असे लिहितो,
  तेव्हा आणि केवळ तेव्हाच, जेव्हा $f \colon A \to B$ असे
  एकास-एक आच्छादन असते.
\end{defn}

\begin{prop}\label{sfr:siz:equ:equinumerosityisequi}
तुल्यबलता हा तुल्यता संबंध आहे.
\end{prop}

\begin{proof}
तुल्यबलता परावर्ती, सममित आणि संक्रमक आहे, हे दाखवावे लागेल. $A, B$
आणि $C$ हे संच असू द्या.

\emph{परावर्तिता.} प्रत्येक $x \in A$ साठी $\Id{A} (x) = x$ असे
असलेले एकरूपता फलन $\Id{A} \colon A \to A$ हे एकास-एक आच्छादन आहे.
म्हणून $\cardeq{A}{A}$.

\emph{सममितता.} $\cardeq{A}{B}$ आहे, म्हणजे $f\colon A \to B$ असे
एकास-एक आच्छादन आहे, असे समजा. $f$ हे एकास-एक व आच्छादक असल्याने त्याचे
व्युत्क्रम $f^{-1}$ अस्तित्वात आहे आणि तेही एकास-एक व आच्छादक आहे. म्हणून
$f^{-1}\colon B \to A$ हे एकास-एक आच्छादन आहे; त्यामुळे
$\cardeq{B}{A}$.

\emph{संक्रमकता.} $\cardeq{A}{B}$ आणि $\cardeq{B}{C}$ आहे, म्हणजे
$f\colon A \to B$ आणि $g\colon B \to
C$ अशी एकास-एक आच्छादने आहेत, असे समजा. मग संयोजन
$\comp{f}{g}\colon A \to C$ हे एकास-एक व आच्छादक आहे; म्हणून
$\cardeq{A}{C}$.
\end{proof}

\begin{prop}
$\cardeq{A}{B}$ असेल, तर $A$ हा गणनीय असतो तेव्हा आणि केवळ
तेव्हाच $B$ तसा असतो.
\end{prop}

\begin{editorial}
पुढील सिद्धतेत \ref{sfr:siz:enm:sec} समाविष्ट असेल, तर
\ref{sfr:siz:enm:defn:enumerable} वापरली जाते; अन्यथा
\ref{sfr:siz:enm-alt:defn:enumerable} वापरली जाते.
\end{editorial}

\begin{proof}
$\cardeq{A}{B}$ आहे, म्हणून $f \colon A
\to B$ असे काही एकास-एक आच्छादन आहे, आणि $A$ हा गणनीय आहे,
असे समजा.

  मग एकतर $A = \emptyset$ असतो, किंवा $g\colon \PosInt \to A$ असे
  आच्छादक फलन असते. $A = \emptyset$ असेल, तर $B = \emptyset$
  सुद्धा असतो (अन्यथा एखादा घटक~$y \in B$ असेल,
  पण $x \in A$ असा कोणताही घटक नसल्याने $f(x) = y$ अशी पूर्वप्रतिमा
  मिळणार नाही). दुसरीकडे, $g\colon \PosInt \to A$ हे आच्छादक
  असेल, तर $\comp{g}{f} \colon \PosInt \to B$ हे आच्छादक आहे.
  हे पाहण्यासाठी $y \in B$ असू द्या. $f$ हे आच्छादक असल्याने
  $f(x) = y$ असा एखादा $x \in A$ घटक असतो. $g$ हे आच्छादक
  असल्याने $g(n) =
  x$ अशी एखादी $n \in \PosInt$ संख्या असते. त्यामुळे
\[(\comp{g}{f})(n) = f(g(n)) = f(x) = y
\]
  आणि म्हणून $\comp{g}{f}$ हे आच्छादक आहे. $\comp{g}{f}$ हे~$B$
  चे प्रगणन आहे; त्यामुळे $B$ हा गणनीय आहे.
\begin{quote}\small\textbf{स्रोतदुरुस्ती.} OLSIZ-006: गोठवलेल्या इंग्रजी सिद्धतेच्या दोन्ही पर्यायी शाखांतील रिक्त-संच प्रकरणात g(x)=y असे लिहिले आहे. त्या शाखेत g अजून अस्तित्वात नसते; A पासून B पर्यंतचे आधी दिलेले द्विएक फलन f वापरूनच B रिक्त असल्याचा निष्कर्ष मिळतो. मराठी मजकुरात दोन्ही ठिकाणी f(x)=y केले आहे.\end{quote}
\begin{quote}\small\textbf{स्रोतदुरुस्ती.} MRSIZ-004: गोठवलेल्या पर्यायी शाखेत एकाच प्रांताला आधी “initial sequence of natural numbers” आणि नंतर “initial segment” म्हटले आहे. आधीच्या औपचारिक व्याख्येशी सुसंगत म्हणून मराठीत दोन्हीकडे “आरंभीचा खंड” वापरला आहे.\end{quote}

$B$ हा गणनीय असेल, तर~$f$ ऐवजी
एकास-एक आच्छादन $f^{-1}\colon B \to A$ घेऊन तोच युक्तिवाद पुन्हा केल्यास
$A$ हा गणनीय आहे, हे मिळते.
\end{proof}

\begin{prob}
$\cardeq{A}{C}$ आणि $\cardeq{B}{D}$ असेल, आणि $A \cap B =
C \cap D = \emptyset$ असेल, तर $\cardeq{A \cup B}{C \cup D}$ आहे,
हे दाखवा.
\end{prob}

\begin{prob}
$A$ हा अनंत आणि गणनीय असेल, तर $\cardeq{A}{\Nat}$ आहे, हे
दाखवा.
\end{prob}









\section{वेगवेगळ्या आकारमानांचे संच आणि कँटर यांचे प्रमेय}\label{sfr:siz:car:sec}

\begin{explain}
दोन संचांचे आकारमान समान असते या कल्पनेचे नेमके विधान आपण दिले आहे.
एक संच दुसऱ्यापेक्षा लहान असतो या कल्पनेचेही नेमके विधान करता येते.
``लहान (किंवा तुल्यबल) असणे'' याच्या व्याख्येत संचांदरम्यानचे
एकास-एक आच्छादन न घेता पहिल्या संचापासून दुसऱ्या संचाकडे जाणारे
एकास-एक फलन घ्यावे लागेल. असे फलन अस्तित्वात असेल, तर पहिल्या
संचाचे आकारमान दुसऱ्या संचाच्या आकारमानापेक्षा लहान किंवा त्याच्याइतके
असते. अंतःप्रज्ञेने, एका संचापासून दुसऱ्या संचाकडे जाणारे
एकास-एक फलन हे फलनाच्या व्याप्तीत प्रांताइतके तरी घटक आहेत
याची खात्री देते, कारण प्रांतातील कोणतेही दोन घटक व्याप्तीतील
एकाच घटक कडे पाठवले जात नाहीत.
\end{explain}

\begin{defn}
$A$ हा $B$ पेक्षा \emph{आकारमानाने मोठा नाही}, हे $\cardle{A}{B}$ असे
लिहितो, तेव्हा आणि केवळ तेव्हाच, जेव्हा $f \colon A \to B$ असे
एकास-एक फलन असते.
\end{defn}

हा संबंध परावर्ती आणि संक्रमक आहे, पण सममित नाही, हे स्पष्ट आहे (हे
सरावासाठी सोडले आहे). एक संच दुसऱ्या संचापेक्षा (काटेकोरपणे) लहान आहे
असे सांगणारी कल्पनाही मांडता येते.

\begin{defn}
$A$ हा $B$ पेक्षा \emph{आकारमानाने लहान} आहे, हे $\cardless{A}{B}$ असे
लिहितो, तेव्हा आणि केवळ तेव्हाच, जेव्हा $f\colon A \to B$ असे
एकास-एक फलन असते, पण $g\colon A
\to B$ असे कोणतेही एकास-एक आच्छादन नसते; म्हणजे, $\cardle{A}{B}$ आणि
$\cardneq{A}{B}$.
\end{defn}

हा संबंध अपरावर्ती आणि संक्रमक आहे, हे स्पष्ट आहे. (हे सरावासाठी
सोडले आहे.) या संकेतनाने संच $A$ हा गणनीय असतो तेव्हा आणि
केवळ तेव्हाच $\cardle{A}{\Nat}$, आणि $A$ हा अगणनीय असतो
तेव्हा आणि केवळ तेव्हाच $\cardless{\Nat}{A}$, असे म्हणता येते. त्यामुळे
%
\ref{sfr:siz:nen-alt:thm:nonenum-pownat}
हे प्रमेय $\cardless{\Nat}{\Pow{\Nat}}$ असे पुन्हा मांडता येते.
खरे तर, हा शेवटचा मुद्दा \emph{पूर्णपणे व्यापक} आहे, हे
Cantor (1892) यांनी सिद्ध केले:

\begin{thm}[कँटर]\label{sfr:siz:car:thm:cantor}
कोणत्याही संच $A$ साठी $\cardless{A}{\Pow{A}}$.
\end{thm}

\begin{proof}
$f(x) = \{x\}$ हे $f \colon A \to \Pow{A}$ असे एकास-एक फलन आहे;
कारण $x \neq y$ असेल, तर विस्तारात्मकतेमुळे $\{x\} \neq \{y\}$ सुद्धा
असते, आणि म्हणून $f(x) \neq f(y)$. त्यामुळे $\cardle{A}{\Pow{A}}$.

\begin{editorial}
\ref{sfr:siz:nen:sec} समाविष्ट असेल, तर आपण संथ सिद्धता देतो; अन्यथा
\ref{sfr:siz:nen-alt:sec} शी जुळणारी अधिक जलद सिद्धता देतो.
\end{editorial}

%
  आता $g\colon A \to
  \Pow{A}$ असे आच्छादक फलन, आणि त्यामुळे एकास-एक व आच्छादक फलनही,
  असू शकत नाही हे दाखवू; म्हणून $\cardneq{A}{\Pow{A}}$. त्यासाठी
  $g\colon A \to \Pow{A}$ असे समजा. $g$ हे सर्वत्र परिभाषित असल्याने
  प्रत्येक $x \in A$ हा $g(x)
  \subseteq A$ अशा उपसंचाकडे पाठवला जातो. $g$ हे आच्छादक असू शकत नाही,
  हे दाखवता येते. त्यासाठी व्याख्येमुळेच~$g$ च्या व्याप्तीत येऊ न
  शकणारा $\overline{A} \subseteq A$ असा उपसंच परिभाषित करू. पुढील संच घ्या:
  \[  \overline{A} = \Setabs{x \in A}{x \notin g(x)}.
  \]
  प्रत्येक $x \in A$ साठी $g(x)$ परिभाषित असल्याने $\overline{A}$ हा
  स्पष्टपणे~$A$ चा सुयोग्य रीतीने परिभाषित उपसंच आहे. पण तो~$g$ च्या
  व्याप्तीत असू शकत नाही. $x \in A$ स्वैर असू द्या; आपण
  $\overline{A} \neq g(x)$ दाखवू. $x \in g(x)$ असेल, तर तो
  $x \notin g(x)$ हा गुणधर्म पूर्ण करत नाही, आणि म्हणून~$\overline{A}$
  च्या व्याख्येवरून $x \notin
  \overline{A}$. $x \in \overline{A}$ असेल, तर त्याने~$\overline{A}$
  चा व्याख्यात्मक गुणधर्म पूर्ण केला पाहिजे; म्हणजे $x \in A$ आणि
  $x \notin
  g(x)$. $x$ स्वैर असल्याने प्रत्येक $x \in A$ साठी $x \in g(x)$ तेव्हा
  आणि केवळ तेव्हाच $x \notin \overline{A}$; म्हणून
  $g(x) \neq \overline{A}$.
  दुसऱ्या शब्दांत, $\overline{A}$ हा~$g$ च्या व्याप्तीत येऊ शकत नाही;
  त्यामुळे~$g$ आच्छादक आहे या गृहीतकाला व्याघात होतो.{}
\begin{quote}\small\textbf{स्रोतदुरुस्ती.} OLSIZ-007: गोठवलेल्या संथ सिद्धतेत x स्वैरपणे A मधून घेतल्यानंतर निष्कर्षाचा आवाका चुकून “प्रत्येक x जो A-bar मध्ये आहे” असा दिला आहे. g(x) आणि A-bar वेगळे आहेत हे प्रत्येक x\ensuremath{\in}A साठी दाखवणे आवश्यक आहे; मराठी सिद्धतेत तो आवाका दुरुस्त केला आहे.\end{quote}
\end{proof}

\begin{explain}

  \ref{sfr:siz:car:thm:cantor} ची सिद्धता आणि
  \ref{sfr:siz:nen:thm:nonenum-pownat} ची सिद्धता यांची तुलना करणे उद्बोधक
  आहे. तेथे~$\PosInt$ च्या उपसंचांची कोणतीही $Z_1$, $Z_2$, \dots{}
  अशी यादी दिली असता, यादीत नसेल असा संख्यांचा संच~$\overline{Z}$
  रचता येतो, हे दाखवले. तो यादीत नसतो याची खात्री यामुळे मिळते की
  प्रत्येक $n \in
  \PosInt$ साठी $n \in Z_n$ तेव्हा आणि केवळ तेव्हाच
  $n \notin \overline{Z}$. त्यामुळे $Z_n$ आणि $\overline{Z}$ यांपैकी
  एकाचा घटक असलेली, पण दुसऱ्याचा नसलेली कोणती तरी संख्या
  नेहमी असते. येथेही तीच कल्पना वापरतो; फरक एवढाच की निर्देशांक~$n$
  आता~$\PosInt$ चे नसून~$A$ चे घटक आहेत. $\overline{A}$ ची
  व्याख्या अशी केली आहे की प्रत्येक $x \in A$ साठी तो~$g(x)$ पेक्षा
  वेगळा असतो, कारण $x \in g(x)$ तेव्हा आणि केवळ तेव्हाच
  $x \notin \overline{A}$. पुन्हा, $g(x)$ आणि $\overline{A}$ यांपैकी
  एकाचा घटक असलेला, पण दुसऱ्याचा नसलेला~$A$ चा कोणता तरी
  घटक नेहमी असतो. म्हणूनच जसा $\overline{Z}$ हा $Z_1$,
  $Z_2$, \dots{} या यादीत येऊ शकत नाही, तसा $\overline{A}$ हा~$g$ च्या
  व्याप्तीत येऊ शकत नाही.


  \ref{sfr:siz:car:thm:cantor} ची सिद्धता आणि
  \ref{sfr:siz:nen-alt:thm:nonenum-pownat} ची सिद्धता यांची तुलना करणे
  उद्बोधक आहे. तेथे~$\Nat$ च्या उपसंचांची कोणतीही $N_0$, $N_1$,
  $N_2$, \dots{} अशी यादी दिली असता, यादीत नसेल असा संख्यांचा संच~$D$
  रचता येतो, हे दाखवले. प्रत्येक $n \in \Nat$ साठी $n \in N_n$ तेव्हा
  आणि केवळ तेव्हाच $n \notin
  D$ असल्यामुळे तो यादीत नसतो याची खात्री मिळते. येथेही तीच कल्पना
  वापरतो; फरक एवढाच की निर्देशांक~$n$ आता~$\Nat$ चे नसून~$A$ चे
  घटक आहेत. $D$ ची व्याख्या अशी केली आहे की प्रत्येक $x
  \in A$ साठी तो~$g(x)$ पेक्षा वेगळा असतो, कारण $x \in g(x)$ तेव्हा
  आणि केवळ तेव्हाच $x \notin D$.

या सिद्धतेची रसेल यांच्या विरोधापत्तीच्या सिद्धतेशीही तुलना करणे उपयुक्त
आहे: \ref{sfr:set:rus:thm:russells-paradox}. खरे तर, कँटर यांच्या
प्रमेयातूनच रसेल यांना स्वतःची विरोधापत्ती सुचली.
\end{explain}

\begin{prob}
  कोणत्याही संच~$A$ साठी $g\colon \Pow{A} \to
  A$ असे एकास-एक फलन असू शकत नाही, हे दाखवा. सूचना: $g\colon
  \Pow{A} \to A$ हे एकास-एक आहे, असे समजा. पुढील संच विचारात घ्या:
  $D = \Setabs{g(B)}{B \subseteq A \text{ आणि
  } g(B) \notin B}$. $x = g(D)$ असे घ्या. $g$ हे एकास-एक आहे हा
  हे तथ्य वापरून व्याघात मिळवा.
\end{prob}









\section{आकारमानाची कल्पना आणि श्रेडर--बर्नस्टाइन प्रमेय}\label{sfr:siz:sb:sec}

\begin{explain}
पुढील विचार अंतर्ज्ञानाला पटणारा आहे: $A$ हा $B$ पेक्षा आकारमानाने मोठा नाही
आणि $B$ हा $A$ पेक्षा आकारमानाने मोठा नाही, तर $A$ आणि $B$ तुल्यबल
आहेत. खरे सांगायचे तर, हा विचार \emph{चुकीचा} असता, तर तुल्यबलतेच्या
आपल्या व्याख्येचा संचांच्या ``आकारमानां''च्या तुलनेशी काही संबंध आहे,
असे म्हणण्याला आपल्याजवळ जवळजवळ आधारच उरला नसता!{} सुदैवाने हा
अंतःप्रज्ञ विचार बरोबर आहे. श्रेडर--बर्नस्टाइन प्रमेय त्याचे समर्थन
करते.
\end{explain}

\begin{thm}[श्रेडर--बर्नस्टाइन]
	\label{sfr:siz:sb:thm:schroder-bernstein}
	$\cardle{A}{B}$ आणि $\cardle{B}{A}$ असेल,
	तर $\cardeq{A}{B}$.
\end{thm}

\begin{explain}
दुसऱ्या शब्दांत, $A$ पासून~$B$ कडे जाणारे एकास-एक फलन आणि $B$
पासून~$A$ कडे जाणारे एकास-एक फलन असेल, तर $A$ पासून~$B$ कडे
जाणारे एकास-एक आच्छादन असते.

पण या निष्कर्षाची सिद्धता खरोखरच बरीच \emph{कठीण} आहे. कँटर यांनी
निष्कर्ष मांडला असला, तरी इतरांनी तो सिद्ध केला.
\footnote{इतिहासाच्या अधिक माहितीसाठी, उदाहरणार्थ,
Potter (2004, pp.~165--6) पाहा.}
%
आत्तासाठी हा निष्कर्ष विश्वासाने स्वीकारावा लागेल.

सुदैवाने श्रेडर--बर्नस्टाइन प्रमेय \emph{बरोबर} आहे आणि म्हणून आपण
परिभाषित केलेले $\cardeq{A}{B}$ आणि $\cardle{A}{B}$ हे संबंध
``आकारमाना''शी निगडित आहेत, असे मानणे समर्थित होते. शिवाय
श्रेडर--बर्नस्टाइन प्रमेय फार \emph{उपयुक्त} आहे. दोन तुल्यबल संचांमध्ये
एकास-एक आच्छादन शोधणे कठीण असू शकते. या प्रमेयामुळे तुलना दोन दिशांत
विभागता येते: दोन्ही दिशांतील एकास-एक फलन स्वतंत्रपणे शोधावी
लागतात.
\end{explain}








\section{प्रगणने आणि गणनीय संच}\label{sfr:siz:enm-alt:sec}

\begin{editorial}
  या विभागात संचसिद्धान्तात रूढ असलेल्या पद्धतीने, $\Nat$ च्या
  (आरंभीच्या खंडां)शी असलेली एकास-एक आच्छादने म्हणून प्रगणनांची
  व्याख्या केली आहे. त्यामुळे \ref{sfr:siz:enm:sec} मधील व्याख्यांशी तिचा
  किंचित विरोध होतो, आणि तेथील सर्व उदाहरणे येथे पुन्हा येतात. हा
  विभाग त्या विभागापेक्षा थोडा अधिक संक्षिप्तही आहे.
\end{editorial}

सांत संच त्यातील घटक केवळ क्रमाने मांडून सांगता येतो. पुढीलप्रमाणे
संचाची व्याख्या करताना आपण हेच करतो:
\[  A = \{a_1, a_2, \ldots, a_n\}.
\]
घटक $a_1$, \dots, $a_n$ हे सर्व भिन्न आहेत असे गृहीत धरल्यास,
यातून $A$ आणि पहिल्या $n$ नैसर्गिक संख्या $0$, \dots, $n-1$ यांच्यात
एकास-एक आच्छादन मिळते. उलट, प्रत्येक सांत संचात सांत संख्येनेच
घटक असल्यामुळे प्रत्येक सांत संच अशी संगती लावून मांडता येतो.
दुसऱ्या शब्दांत, $A$ सांत असेल, तर $A$ आणि $\{0, \dots, n-1\}$
यांच्यात एकास-एक आच्छादन असते; येथे~$A$ मध्ये $n$ इतके घटक
असतात.

विशिष्ट प्रकारचे अनंत संच मान्य केल्यास, काही अनंत संचांचेही प्रगणन
करता येईल. अनंत संच आणि संपूर्ण~$\Nat$ यांच्यातील एकास-एक आच्छादन त्या
संचाचे प्रगणन करते, असे म्हणून हे नेमकेपणाने मांडता येते.

\begin{defn}[प्रगणनाची संचसैद्धान्तिक व्याख्या]
संच $A$ चे \emph{प्रगणन} म्हणजे असे एकास-एक आच्छादन की ज्याची व्याप्ती
$A$ असते आणि ज्याचा प्रांत नैसर्गिक संख्यांचा एखादा आरंभीचा खंड $\{0,
1, \ldots, n\}$ {किंवा} नैसर्गिक संख्यांचा संपूर्ण संच~$\Nat$ असतो.
\end{defn}

\begin{explain}
\emph{प्रगणन} या शब्दाच्या अशा वापरामागे सहज समजणारा आधार आहे.
संच $A$ चे प्रगणन केले आहे असे म्हणणे म्हणजे संच~$A$ मधील घटक मोजत
जाण्यास अनुमती देणारे एकास-एक आच्छादन $f$ आहे असे म्हणणे. $0$ वा घटक
$f(0)$, पहिला $f(1)$, \ldots आणि $n$ वा $f(n)$,~\ldots असतो.
\footnote{होय, आपण $0$ पासून मोजतो. अर्थात~$1$ पासूनही सुरुवात करता
येईल. त्यामुळे मोठा फरक पडणार नाही; फक्त~$\Nat$ ऐवजी~$\PosInt$
घ्यावा लागेल.} पुढील व्याख्या जोडल्यावर या मागचे कारण आणखी स्पष्ट होते:
\end{explain}

\begin{defn}
  \label{sfr:siz:enm-alt:defn:enumerable}
  संच~$A$ हा गणनीय असतो तेव्हा आणि केवळ तेव्हाच, जेव्हा
  $A = \emptyset$ असते किंवा~$A$ चे प्रगणन असते. $A$ हा
  अगणनीय आहे असे म्हणतो तेव्हा आणि केवळ तेव्हाच, जेव्हा
  $A$ हा गणनीय नसतो.
\end{defn}

\begin{explain}
म्हणून संच रिक्त असेल किंवा प्रगणन वापरून त्यातील घटक मोजता
येत असतील तेव्हा आणि केवळ तेव्हाच तो गणनीय असतो.
\end{explain}

\begin{ex}
नैसर्गिक संख्यांचे प्रगणन करणारे फलन म्हणजे फक्त $\Id{\Nat} \colon
\Nat \to \Nat$ हे $\Id{\Nat}(n) = n$ ने दिलेले एकरूपता फलन.
\emph{धन} नैसर्गिक संख्या $\Nat^+ = \Nat \setminus \{0\}$ यांचे
प्रगणन करणारे फलन $g(n) = n + 1$ आहे; म्हणजेच ते उत्तरवर्ती फलन आहे.
\end{ex}

\begin{prob}
संच $A$ हा गणनीय असतो तेव्हा आणि केवळ तेव्हाच $A =
\emptyset$ असते किंवा $f\colon \Nat \to A$ असे आच्छादन
असते, हे दाखवा. $A$ हा गणनीय असतो तेव्हा आणि केवळ तेव्हाच
$g\colon A \to \Nat$ असे एकास-एक फलन असते, हेही दाखवा.
\end{prob}

\begin{ex}
पुढीलप्रमाणे दिलेली $f\colon \Nat \to \Nat$ आणि $g \colon \Nat
\to \Nat$ ही फलने
\begin{align*}
  f(n) & = 2n \text{ आणि}\\
  g(n) & = 2n+1
\end{align*}
अनुक्रमे सम नैसर्गिक संख्यांचे आणि विषम नैसर्गिक संख्यांचे प्रगणन
करतात. पण यांपैकी कोणतेही फलन आच्छादक नाही; म्हणून कोणतेही
$\Nat$ चे प्रगणन नाही.
\end{ex}

\begin{prob}
वर्गसंख्या $1$, $4$, $9$, $16$, \dots{} यांचे प्रगणन परिभाषित करा.
\end{prob}

\begin{ex}
$\lceil x \rceil$ हे $x$ ला त्याच्यापेक्षा कमी नसलेल्या सर्वांत जवळच्या
पूर्णांकाकडे नेणारे \emph{ऊर्ध्व पूर्णांक} फलन असू द्या. मग पुढीलप्रमाणे
दिलेले फलन $f \colon \Nat \to \Int$:
\[  f(n) = (-1)^{n} \left\lceil\tfrac{n}{2}\right\rceil
\]
पूर्णांकांचा संच~$\Int$ पुढीलप्रमाणे प्रगणित करते:
\[\begin{array}{c c c c c c c c}
f(0) & f(1) & f(2) & f(3) & f(4) & f(5) & f(6) & \dots \\ \\
\big\lceil \tfrac{0}{2} \big\rceil & -\big\lceil \tfrac{1}{2}\big\rceil &  \big\lceil \tfrac{2}{2} \big\rceil & -\big\lceil \tfrac{3}{2} \big\rceil & \big\lceil \tfrac{4}{2} \big\rceil  & -\big\lceil \tfrac{5}{2}\big\rceil & \big\lceil \tfrac{6}{2} \big\rceil & \dots \\ \\
0 & -1 & 1 & -2 & 2 & -3 & 3& \dots
\end{array}
\]
धन आणि ऋण पूर्णांकांमध्ये आलटूनपालटून ``उड्या मारून'' $f$ हे $\Int$
ची मूल्ये कशी निर्माण करते ते पाहा. $f$ ची व्याख्या पुढीलप्रमाणे
प्रकरणे पाडून केली आहे असेही मानता येते:
\[f(n) = \begin{cases}
  \frac{n}{2} & \text{जर $n$ सम असेल}\\
  -\frac{n+1}{2} & \text{जर $n$ विषम असेल}
  \end{cases}
\]
\end{ex}

\begin{prob}
$A$ आणि $B$ हे गणनीय असतील, तर $A \cup B$ सुद्धा गणनीय
आहे, हे दाखवा.
\end{prob}

\begin{prob}
$A_1$, $A_2$, \dots, $A_n$ हे सर्व गणनीय असतील, तर
$A_1 \cup \dots \cup A_n$ सुद्धा गणनीय आहे, हे $n$ वरील गणिती विगमनाने
दाखवा.
\end{prob}









\section{अगणनीय संच}\label{sfr:siz:nen-alt:sec}

\begin{editorial}
  या विभागात \ref{sfr:siz:enm-alt:sec} मधील व्याख्या वापरून
  $\Bin^\omega$ आणि $\Pow{\Nat}$ यांची अगणनीयता सिद्ध केली आहे;
  म्हणजे $\PosInt$ पासून आच्छादन मागण्याऐवजी~$\Nat$ शी एकास-एक
  आच्छादन आवश्यक मानले आहे.
\end{editorial}

\begin{explain}
नैसर्गिक संख्यांचा संच $\Nat$ हा अनंत आहे. तो सहजपणे
गणनीय देखील आहे. परंतु
\emph{अगणनीय} संच, म्हणजे गणनीय नसलेले संच,
अस्तित्वात आहेत ही लक्षणीय गोष्ट आहे
(पाहा \ref{sfr:siz:enm-alt:defn:enumerable}).

हे आश्चर्यकारक वाटू शकते. कारण $A$ हा अगणनीय आहे असे
म्हणणे म्हणजे \emph{कोणतेही} एकास-एक आच्छादन $f \colon \Nat \to A$
नाही असे म्हणणे; म्हणजे~$\Nat$ मधील अनंतपणे अनेक घटक~$A$
मध्ये पाठवणारे कोणतेही फलन संपूर्ण~$A$ व्यापत नाही. म्हणून $A$ हा
अगणनीय असेल, तर नैसर्गिक संख्यांपेक्षा~$A$ मध्ये ``अधिक''
घटक असतात.

संच अगणनीय आहे हे सिद्ध करण्यासाठी अनुरूप एकास-एक आच्छादन
अस्तित्वात असू शकत नाही हे दाखवावे लागते. याचा सर्वोत्तम मार्ग म्हणजे
~$A$ मधील घटक प्रगणित करण्याचा प्रत्येक प्रयत्न किमान एक
घटक वगळतो हे दाखवणे; यावरून कोणतेही फलन $f\colon \Nat \to A$
हे आच्छादक नाही असे दिसते. हे प्रस्थापित करण्याची एक सर्वसाधारण
युक्ती म्हणजे कँटर यांची \emph{कर्णरेषीय पद्धत} वापरणे. $A$ मधील
घटक $x_1$, $x_2$, \dots{} असे मांडले असता, रचनेमुळे त्या
यादीत असूच शकणार नाही असा~$A$ चा आणखी एक घटक आपण बनवतो.

हे सर्व उदाहरणाने उत्तम समजते. म्हणून आपले पहिले उदाहरण म्हणजे $0$
आणि $1$ यांच्या सर्व अनंत चिन्हमालांचा संच~$\Bin^\omega$. (`$\Bin$'
हे द्विमान दर्शवते आणि त्याचा अर्थ फक्त दोन घटकांचा संच $\{0,1\}$
असा घेता येतो.)
तरी सध्या या किंचित शिथिल मांडणीमुळे गोंधळ होऊ नये.
\end{explain}

\begin{thm}
\label{sfr:siz:nen-alt:thm:nonenum-bin-omega}
$\Bin^\omega$~हा अगणनीय आहे.
\end{thm}

\begin{proof}
$\Bin^\omega$ च्या एखाद्या उपसंचाचे कोणतेही प्रगणन घ्या. म्हणजे
$s_{0}$, $s_{1}$, $s_{2}$, \dots{} अशी यादी आपल्याकडे आहे, जेथे
प्रत्येक $s_n$ ही $0$ आणि~$1$ यांची अनंत चिन्हमाला आहे. या यादीतील
$n$ व्या चिन्हमालेतील $m$ वा अंक $s_n(m)$ असू द्या.
\begin{quote}\small\textbf{स्रोतदुरुस्ती.} OLSIZ-008: गोठवलेल्या इंग्रजी वर्णनात s\_n(m) ला m व्या चिन्हमालेतील n वा अंक म्हटले आहे; लगेचचे वाक्य आणि सारणी मात्र तो n व्या ओळीतील, म्हणजे n व्या चिन्हमालेतील, m वा अंक असल्याचे निश्चित करतात. मराठी मजकुरात ओळ-स्तंभाशी सुसंगत निर्देशांकक्रम वापरला आहे.\end{quote}
मग यादीकडे सरणी म्हणून पाहता येते, ज्यात $s_n(m)$ हा $n$ व्या ओळीत
आणि $m$ व्या स्तंभात ठेवलेला आहे:
\[\begin{array}{c|c|c|c|c|c}
& 0 & 1 & 2 & 3 & \dots \\\hline
0 & \mathbf{s_{0}(0)} & s_{0}(1) & s_{0}(2) & s_0(3) & \dots \\\hline
1 & s_{1}(0)& \mathbf{s_{1}(1)} & s_1(2) & s_1(3) & \dots \\\hline
2 & s_{2}(0)& s_{2}(1) & \mathbf{s_2(2)} & s_2(3) & \dots \\\hline
3 & s_{3}(0)& s_{3}(1) & s_3(2) & \mathbf{s_3(3)} & \dots \\\hline
\vdots & \vdots & \vdots & \vdots & \vdots & \mathbf{\ddots}
\end{array}
\]
आता या यादीत नसलेली $0$ आणि $1$ यांची $d$ ही अनंत चिन्हमाला आपण
बनवू. तिची प्रत्येक नोंद सांगून, म्हणजे सर्व $n \in \Nat$ साठी
$d(n)$ सांगून, हे करू. सहज कल्पना अशी: वरील सरणीचा कर्ण खाली वाचायचा
(त्यामुळे ``कर्णरेषीय पद्धत'' हे नाव) आणि प्रत्येक $1$ चा $0$, तर
प्रत्येक $0$ चा~$1$ करायचा.
\begin{quote}\small\textbf{स्रोतदुरुस्ती.} OLSIZ-009: गोठवलेल्या इंग्रजी सूचनेत “प्रत्येक 1 चा 0” हेच परिवर्तन दोनदा आले आहे. लगेचचे प्रकरणनिहाय सूत्र परस्परपूरक दोन्ही बदल निश्चित करते: 1 चा 0 आणि 0 चा 1. मराठी मजकुरात हे दोन्ही बदल स्पष्ट केले आहेत.\end{quote}
अधिक अमूर्तपणे, कर्णावरील $n$ वा घटक, $s_n(n)$, हा $1$ किंवा
$0$ आहे त्यानुसार $d(n)$ हा $0$ किंवा $1$ असावा अशी व्याख्या करतो;
म्हणजे:
\[d(n) =
\begin{cases}
1 & \text{जर $s_{n}(n) = 0$ असेल}\\
0 & \text{जर $s_{n}(n) = 1$ असेल}
\end{cases}
\]
स्पष्टच, $d \in \Bin^\omega$, कारण ती $0$ आणि $1$ यांची अनंत
चिन्हमाला आहे. पण आपण $d$ अशी बनवली की कोणत्याही $n \in \Nat$
साठी $d(n) \neq s_n(n)$. म्हणजे $d$ ही $s_n$ पेक्षा तिच्या $n$ व्या
नोंदीत वेगळी आहे. म्हणून कोणत्याही $n\in \Nat$ साठी $d \neq s_n$.
त्यामुळे $d$ ही $s_0$, $s_1$, $s_2$, \dots{} या यादीत असू शकत नाही.

$\Bin^\omega$ च्या एखाद्या उपसंचाचे स्वैर प्रगणन दिले असता ते
$\Bin^\omega$ चा कोणता तरी घटक वगळेल हे आपण दाखवले आहे.
म्हणून $\Bin^\omega$ या संचाचे प्रगणन नाही; म्हणजे $\Bin^\omega$ हा
अगणनीय आहे.
\end{proof}

\begin{explain}
या सिद्धतापद्धतीला ``कर्णीकरण'' म्हणतात, कारण ती~$d$ ची व्याख्या
करण्यासाठी सरणीचा कर्ण वापरते. मात्र कर्णीकरणात सरणी असणे आवश्यक नाही.
प्रत्यक्षात, सरणी किंवा प्रत्यक्ष कर्ण नसतानाही हिच्याशी मिळतीजुळती कल्पना
वापरून एखादा संच अगणनीय आहे हे दाखवता येते. पुढील फलित हे
कसे करता येते ते दाखवते.
\end{explain}

\begin{thm}
\label{sfr:siz:nen-alt:thm:nonenum-pownat}
$\Pow{\Nat}$ हा गणनीय नाही.
\end{thm}

\begin{proof}
~$\Nat$ च्या उपसंचांची प्रत्येक यादी $\Nat$ चा कोणता तरी उपसंच वगळते
हे दाखवून आपण त्याच पद्धतीने पुढे जाऊ. म्हणून $\Nat$ च्या उपसंचांची
$N_0, N_1, N_2, \ldots$ अशी एखादी यादी आहे असे समजा. $D$ या संचाची
व्याख्या पुढीलप्रमाणे करतो: $n \in D$ तेव्हा आणि केवळ तेव्हाच, जेव्हा
$n \notin N_{n}$:
\[D = \Setabs{n \in \Nat}{n \notin N_n}
\]
स्पष्टच $D\subseteq \Nat$. पण $D$ हा यादीत असू शकत नाही. कारण
रचनेनुसार $n \in D$ तेव्हा आणि केवळ तेव्हाच, जेव्हा $n\notin N_n$;
म्हणून कोणत्याही $n \in \Nat$ साठी $D \neq N_n$.
\end{proof}

\begin{explain}
आधीच्या सिद्धतेत कर्णाचा उल्लेख नव्हता. तरी पुढीलप्रमाणे चित्र डोळ्यांसमोर
आणल्यास त्यात कर्ण आहे असे मानता येते: $N_0$, $N_1$, \dots{} हे संच
एका सरणीत लिहिल्याची कल्पना करा; $m \in N_n$ असेल तेव्हा आणि केवळ
तेव्हाच $m$ व्या स्तंभात $m$ लिहून $N_n$ हा $n$ व्या ओळीत लिहायचा.
उदाहरणार्थ, त्या यादीतील पहिले चार संच $\{0,1,2,\dots\}$,
$\{1, 3, 5, \dots\}$, $\{0,1,4\}$ आणि $\{2,3,4,\dots\}$ आहेत असे
समजा; मग आपल्या सरणीची सुरुवात अशी असेल:
\[\begin{array}{r@{}rrrrrrr}
  N_0 = \{ & \mathbf{0}, & 1, & 2, & & & & \dots\}\\
  N_1 = \{ &  & \mathbf{1}, &  & 3, &  & 5, & \dots\}\\
  N_2 = \{ & 0, & 1, &  &  & 4\phantom{,} &  & \}\\
  N_3 = \{ &  &  & 2, & \mathbf{3}, & 4, & & \dots\}\\
  &\vdots & & & & & & \ddots\phantom{\}}
  \end{array}
\]
मग कर्णाच्या खाली जात, $n$ कर्णावर \emph{नसेल} तेव्हा आणि केवळ
तेव्हाच $n \in D$ असे ठेवून $D$ हा संच मिळतो. म्हणून वरील उदाहरणात
आपण $0$ आणि $1$ वगळू,~$2$ समाविष्ट करू,~$3$ वगळू, इत्यादी.
\end{explain}

\begin{prob}
सर्व फलने $f \colon \Nat \to \Nat$ यांचा संच स्पष्ट कर्णरेषीय
युक्तिवादाने अगणनीय आहे हे दाखवा. म्हणजे $f_1$, $f_2$,
\dots{} ही फलनांची यादी असून प्रत्येक $f_i\colon \Nat \to \Nat$
असेल, तर या यादीत नसलेले कोणते तरी $g \colon \Nat \to \Nat$ असते,
हे दाखवा.
\end{prob}









\section{न्यूनीकरण}\label{sfr:siz:red-alt:sec}

\begin{editorial}
  या विभागात \ref{sfr:siz:nen-alt:sec} मधील फलितांशी जुळणारी अगणनीयता
  न्यूनीकरणाने सिद्ध केली आहे. \ref{sfr:siz:nen:sec} मधील फलितांशी जुळणारी
  किंचित अधिक विस्तृत पर्यायी आवृत्ती \ref{sfr:siz:red:sec} मध्ये दिली आहे.
\end{editorial}

कर्णीकरणाच्या युक्तिवादाने $\Bin^\omega$ हा अगणनीय आहे,
हे आपण सिद्ध केले. तसाच कर्णीकरणाचा युक्तिवाद वापरून $\Pow{\Nat}$ हा
अगणनीय आहे हेही दाखवले. पण $\Pow{\Nat}$ हा
अगणनीय असल्याचे सिद्ध करण्याचा आणखी एक मार्ग आहे:
\emph{$\Pow{\Nat}$ हा गणनीय असेल, तर $\Bin^\omega$ हाही
गणनीय असतो} हे दाखवा. $\Bin^\omega$ हा अगणनीय
आहे हे आपल्याला माहीत असल्यामुळे $\Pow{\Nat}$ हाही तसाच आहे असा
निष्कर्ष निघेल.

याला एक समस्या दुसऱ्या समस्येकडे \emph{न्यूनीत करणे} म्हणतात. या
बाबतीत $\Bin^\omega$ चे प्रगणन करण्याची समस्या आपण $\Pow{\Nat}$ चे
प्रगणन करण्याच्या समस्येकडे न्यूनीत करतो. नंतरच्या समस्येचे उत्तर---
$\Pow{\Nat}$ चे प्रगणन---पहिल्या समस्येचे उत्तर---$\Bin^\omega$ चे
प्रगणन---देईल.

संच~$B$ च्या प्रगणनाची समस्या संच~$A$ च्या प्रगणनाच्या समस्येकडे
न्यूनीत करण्यासाठी~$A$ चे प्रगणन~$B$ च्या प्रगणनात रूपांतरित करण्याची
पद्धत आपण देतो. हे करण्याचा सर्वांत सोपा मार्ग म्हणजे
$f\colon A \to B$ असे आच्छादन परिभाषित करणे. $x_1$, $x_2$,
\dots{} हे~$A$ चे प्रगणन असेल, तर $f(x_1)$, $f(x_2)$, \dots{} हे~$B$
चे प्रगणन करील. आपल्या बाबतीत आपण
$f\colon \Pow{\Nat} \to \Bin^\omega$ असे आच्छादन शोधत आहोत.

\begin{prob}
$g\colon B \to A$ असे एकास-एक फलन असेल आणि $B$ हा
अगणनीय असेल, तर $A$ हाही तसाच असतो, हे दाखवा. $g$ वापरून
~$A$ चे प्रगणन~$B$ च्या प्रगणनात कसे रूपांतरित करता येते हे दाखवून
सिद्धता द्या.
\end{prob}

\begin{proof}[न्यूनीकरणाने {\ref{sfr:siz:nen-alt:thm:nonenum-pownat}} ची सिद्धता]
न्यूनीकरणासाठी, $\Pow{\Nat}$ हा गणनीय आहे आणि त्यामुळे त्याचे
$N_{1}$, $N_{2}$, $N_{3}$, \dots{} असे प्रगणन आहे, असे समजा.

$f \colon \Pow{\Nat} \to \Bin^\omega$ हे फलन पुढीलप्रमाणे परिभाषित
करा: $f(N)$ ही अशी चिन्हमाला~$s$ असू द्या की $s(n) = 1$ तेव्हा आणि
केवळ तेव्हाच, जेव्हा $n \in N$; अन्यथा $s(n) = 0$.
\begin{quote}\small\textbf{स्रोतदुरुस्ती.} OLSIZ-010: गोठवलेल्या इंग्रजी व्याख्येत f(N) ला s\_k असे म्हटले आहे, पण k बांधलेला किंवा N वरून ठरवलेला नाही. घटकनिहाय समीकरण N ची अभिलक्षण-चिन्हमाला एकमेव ठरवते; मराठी व्याख्येत एकच बद्ध निर्गत नाव s आणि त्याचे घटक s(n) सातत्याने वापरले आहेत.\end{quote}

याने फलन स्पष्टपणे परिभाषित होते, कारण $N \subseteq \Nat$ असताना
कोणताही $n \in \Nat$ हा $N$ चा घटक असतो किंवा नसतो.
उदाहरणार्थ, सम नैसर्गिक संख्यांचा संच
$2\Nat = \Setabs{2n}{n \in \Nat} = \{0,2, 4, 6, \dots\}$ हा
$1010101\dots$ या चिन्हमालेकडे; $\emptyset$ हा $0000\dots$ कडे; आणि
$\Nat$ हा $1111\dots$ कडे पाठवला जातो.

हे फलन आच्छादक देखील आहे: $0$ आणि $1$ यांची प्रत्येक चिन्हमाला
नैसर्गिक संख्यांच्या कोणत्या तरी संचाशी अनुरूप असते; तो संच म्हणजे
चिन्हमालेत ज्या स्थानांवर~$1$ आहे त्यांना अनुरूप नैसर्गिक संख्या
सदस्य म्हणून असलेला संच. अधिक नेमकेपणाने, $s \in \Bin^\omega$ असेल,
तर पुढीलप्रमाणे $N \subseteq \Nat$ परिभाषित करा:
\[N = \Setabs{n \in \Nat}{s(n) = 1}
\]
मग~$f$ ची व्याख्या पाहून $f(N) = s$ हे तपासता येते.

आता पुढील यादी विचारात घ्या:
\[f(N_1), f(N_2), f(N_3), \dots
\]
$f$ हे आच्छादक असल्यामुळे $\Bin^\omega$ चा प्रत्येक सदस्य
हा~$f$ चे कोणत्या तरी फलसाधकावरील मूल्य असला पाहिजे आणि म्हणून यादीत
आला पाहिजे. त्यामुळे ही यादी संपूर्ण~$\Bin^\omega$ चे प्रगणन करते.

म्हणून $\Pow{\Nat}$ हा गणनीय असता, तर $\Bin^\omega$ हाही
गणनीय असता. पण $\Bin^\omega$ हा अगणनीय आहे
(\ref{sfr:siz:nen-alt:thm:nonenum-bin-omega}). त्यामुळे $\Pow{\Nat}$ हा
अगणनीय आहे.
\end{proof}


\begin{prob}\label{sfr:siz:red:prob:nat-nat-alt}
  $f\colon \Nat \to \Nat$ अशा सर्व फलनांचा संच~$X$ हा
  न्यूनीकरणाच्या युक्तिवादाने अगणनीय आहे, हे दाखवा.
  (सूचना: $X$ पासून~$\Bin^\omega$ कडे जाणारे आच्छादक फलन द्या.)
\end{prob}

\begin{prob}
नैसर्गिक संख्यांच्या जोड्यांच्या \emph{सर्व संचांचा} संच, म्हणजे
$\Pow{\Nat \times \Nat}$, न्यूनीकरणाच्या युक्तिवादाने
अगणनीय आहे, हे दाखवा.
\end{prob}

\begin{prob}
नैसर्गिक संख्यांच्या अनंत अनुक्रमांचा संच $\Nat^\omega$ हा
न्यूनीकरणाच्या युक्तिवादाने अगणनीय आहे, हे दाखवा.
\end{prob}


\begin{prob}
$S$ हा $\Nat$ पासून $\{0,1\}$ या संचाकडे जाणारी सर्व
आच्छादक फलने असलेला संच असू द्या; म्हणजे $S$ मध्ये
$f \colon \Nat \to \Bin$ अशी सर्व आच्छादक फलने आहेत. $S$ हा
अगणनीय आहे, हे दाखवा.
\end{prob}

\begin{prob}
सर्व वास्तव संख्यांचा संच~$\Real$ हा अगणनीय आहे, हे दाखवा.
\end{prob}



